% concatenated from Isomorphism_of_GBD.tex
\documentclass[11pt, twoside]{amsart}
\usepackage[cp1251]{inputenc}
\usepackage{color}
\usepackage{amsthm,amssymb,latexsym}
\usepackage[centertags]{amsmath}
%\usepackage{hyperref,url}
%\usepackage{graphicx}
\usepackage[T2A]{fontenc}
\usepackage{graphs} 
\usepackage{MnSymbol}
\usepackage{bbm}
\usepackage{centernot}
\usepackage{mathrsfs}
\usepackage{amsthm}
\usepackage{amsfonts}
\usepackage{textcomp}
\usepackage{newlfont}
\usepackage{enumitem}   


%-------------pdflatex or latex--------------------------------------
\newif\ifPDF
\ifx\pdfoutput\undefined\PDFfalse
\else \ifnum \pdfoutput > 0 \PDFtrue
        \else \PDFfalse
        \fi
\fi

%------------



%------------packages for pdflatex------------------------------

\ifPDF
  \usepackage[pdftex]{xcolor, graphicx}
  \usepackage[colorlinks=true,linkcolor=blue, citecolor=blue]{hyperref}%

 % \usepackage[pdftex, bookmarks, colorlinks]{hyperref}
  %\hypersetup{
   %  linkcolor = blue,
  %   filecolor = magenta,      
  %   urlcolor = cyan,
  %   pdftitle = {},
   %  bookmarks = true,
 %    pdfpagemode = FullScreen,
%}



 %  \hypersetup{colorlinks=false}

%------------\packages for latex---------------------------------

\else
  \usepackage{color}
  \usepackage[dvips]{graphicx}
  \usepackage[dvips]{hyperref}
\fi


\usepackage{tkz-graph}
\tikzset{EdgeStyle/.style = {->}}
\tikzset{LabelStyle/.style= {fill=yellow}}
\usetikzlibrary{shapes,snakes,calendar,matrix,backgrounds,folding}

\usepackage{bbm}

%-------------------------------------PAGE LAYOUT---------------------------------------------------
%\textheight = 575pt
%\textwidth = 425pt
\usepackage[top=1in, bottom=1.25in, left=1.25in, right=1.25in]
{geometry}
%\textwidth=16cm
%\usepackage[scale=0.75]{geometry}

 
%	\addtolength{\topmargin}{0.3in}
	%\addtolength{\textheight}{1.75in}

\renewcommand{\baselinestretch}{1.1}

%\tolerance =10000
%\hbadness =10000
%\linespread{1.1}
%\textwidth=32cc
%\textheight=225truemm
%\topmargin=-1cm
%\oddsidemargin=1cm
%\evensidemargin=1cm

\theoremstyle{definition}
\newtheorem{definition}{Definition}[section]
\theoremstyle{remark}
\newtheorem{example}[definition]{Example}
\newtheorem{ex}[definition]{Exercise}
\newtheorem*{probl}{Open problem}
\newtheorem{remark}[definition]{Remark}
\newtheorem{quest}[definition]{Question}

\theoremstyle{plain}
\newtheorem{thm}[definition]{Theorem}
\newtheorem{theorem}[definition]{Theorem}
\newtheorem{prop}[definition]{Proposition}
\newtheorem{lemma}[definition]{Lemma}
\newtheorem{fact}[definition]{Fact}
\newtheorem{corol}[definition]{Corollary}
\newtheorem{claim}[definition]{Claim}


\def\rank{{\rm rank}}
\newcommand{\tcb}{\textcolor{blue}}
\newcommand{\tcr}{\textcolor{red}}
\newcommand{\tcg}{\textcolor{green}}

\newcommand{\stirlingii}{\genfrac{\{}{\}}{0pt}{}}

\def\wt{\widetilde}
\def\ol{\overline}
\def\ov{\overline}
\def\tl{\widetilde}
\def\wh{\widehat}
\def\ho{\widehat{\overline}}
\newcommand{\N}{\mathbb N}
\newcommand{\Z}{\mathbb Z}
\newcommand{\E}{\mathcal E}
\newcommand{\B}{\mathcal B}
\newcommand{\om}{\omega}
\newcommand{\ent}{f^{(n)}_{vw}}

\newcommand{\be}{\begin{equation}}
\newcommand{\ee}{\end{equation}}
\newcommand{\ba}{\begin{aligned}}
\newcommand{\ea}{\end{aligned}}
\newcommand{\mc}{\mathcal}


\newcommand{\vp}{\varphi}
\newcommand{\e}{\varepsilon}

%------------------------------------------------------

\newcommand{\La}{\Lambda}
\newcommand{\Om}{\Omega}
\newcommand{\al}{\alpha}
\newcommand{\G}{\Gamma}
\newcommand{\g}{\gamma}
\newcommand{\T}{\theta}
\newcommand{\De}{\Delta}
\newcommand{\de}{\delta}
\newcommand{\s}{\sigma}
\newcommand{\A}{{\cal A}}
\newcommand{\M}{{\cal M}}
\newcommand{\bs}{(X,{\cal B})}
\newcommand{\Aut}{Aut(X,{\cal B})}
%\newcommand{\h}{Homeo(\Om)}
\newcommand{\pos}{{\mathbb R}^*_+}
\newcommand{\R}{{\mathbb R}}
%\newcommand{\la}{\lambda}
\newcommand{\ap}{{\cal A}p}
\newcommand{\per}{{\cal P}er}
\newcommand{\inc}{{\cal I}nc}

\numberwithin{equation}{section}
\newcommand{\ignore}[1]{}

\begin{document}


\title{Isomoprhism of generalized Bratteli diagrams}


\author{Olena Karpel}
%    Address of record for the research reported here
\address{AGH University of Krakow, Faculty of Applied Mathematics, al. Adama Mickiewicza~30, 30-059 Krak\'ow, Poland \&
B. Verkin Institute for Low Temperature Physics and Engineering,
47~Nauky Ave., Kharkiv, 61103, Ukraine}

\email{okarpel@agh.edu.pl}


\subjclass[2020]{37A05, 37B05, 37A40, 54H05, 05C60}

\keywords{Borel dynamical systems, Bratteli-Vershik model, tail-invariant measures, infinite ergodic theory.}


\date{}

\begin{abstract}
%We study the notion of isomorphism for generalized Bratteli diagrams. We consider properties of Bratteli diagrams that are preserved under isomorphism. We show that every generalized Bratteli diagram is isomorphic to an irreducible generalized Bratteli diagram. We introduce the notion of a completely irreducible generalized Bratteli diagram, defined as a diagram that can be isomorphic only to irreducible generalized Bratteli diagrams. We show how this notion is connected to the topological properties of the path space of a generalized Bratteli diagram and its tail equivalence relation. We study in detail several classes of generalized Bratteli diagrams that illustrate the results.


We study the notion of isomorphism for generalized Bratteli diagrams and investigate properties preserved under isomorphism. We show that every generalized Bratteli diagram is isomorphic to an irreducible generalized Bratteli diagram. We introduce the notion of a completely irreducible generalized Bratteli diagram, namely, a diagram that is isomorphic only to irreducible generalized Bratteli diagrams. We establish connections between this notion and the topological properties of the path space of a generalized Bratteli diagram and its tail equivalence relation. We examine in detail several classes of generalized Bratteli diagrams that illustrate these results.
%We apply the results to find new measures on generalized Bratteli diagrams of bounded size.

%see for standard bd reducible and irreducible

\end{abstract}


\maketitle
\tableofcontents


\section{Introduction}\label{intro}


Bratteli diagrams are infinite graded graphs introduced by O. Bratteli in \cite{Bratteli1972} for the study of approximately finite-dimensional $C^*$-algebras. Dynamical systems were associated with Bratteli diagrams by A. Vershik in \cite{Vershik_1981, Vershik_1982}. It turned out that Bratteli diagrams and the dynamical systems arising on their path spaces provide a powerful framework for modeling a broad class of equivalence relations and dynamical systems in measurable, Cantor and Borel dynamics. %(see e.g. \cite{HermanPutnamSkau1992, Medynets_2006, DownarowiczKarpel_2019,Shimomura2020}). 
%For instance, 
One of the fundamental results states that any minimal homeomorphism of a Cantor space can be modelled as a Vershik map acting on the path space of an associated simple  Bratteli diagram \cite{HermanPutnamSkau1992}. 
The combinatorial structure of Bratteli diagrams provides tools for studying invariant measures and orbit structure for the corresponding dynamical systems and equivalence relations. There is vast literature on Bratteli diagrams (see, e.g., the books and surveys \cite{Durand2010, BezuglyiKarpel2016, DownarowiczKarpel2018, Putnam2018, BezuglyiKarpel_2020, DurandPerrin2022} and references therein). 

This paper focuses on the study of generalized Bratteli diagrams, a relatively new notion introduced in \cite{BezuglyiDooleyKwiatkowski2006}. 
In that work, it was shown that any aperiodic Borel automorphism of a standard Borel space 
%is Borel isomorphic to 
can be modelled as a Vershik map acting on the path space of an associated generalized Bratteli diagram. The orbits of the constructed Vershik map coincide with the orbits of the so-called tail equivalence relation associated with the diagram.
While introducing the Vershik map requires introducing partial order on the set of edges of a Bratteli diagram, the tail equivalence relation does not require the order and can be introduced for every (standard or generalized) Bratteli diagram.
Moreover, with every generalized Bratteli diagram one can associate a Borel automorphism whose orbits coincide with those of the associated tail equivalence relation \cite{Kechris2024}. 
It also follows from the results of \cite{BezuglyiDooleyKwiatkowski2006} that every aperiodic hyperfinite countable Borel equivalence relation can be represented as a tail equivalence relation on the path space of a generalized Bratteli diagram.
The structure of a Bratteli diagram provides tools for describing invariants for Borel isomorphism of such equivalence relations, and for orbit equivalence of the corresponding Borel automorphisms. One of such invariants is the number of  ergodic invariant probability measures for the tail equivalence relation
\cite{Dougherty_Jackson_Kechris1994, Kechris2024, BezuglyiKarpelKwiatkowski2019, BezuglyiKarpelKwiatkowskiWata2024}.
In this paper, we entirely focus on the tail equivalence relation, and do not consider the corresponding Vershik map. 


The systematic study of generalized Bratteli diagrams has started only recently (see, e.g.,  \cite{BezuglyiJorgensenSanadhya2024, BezuglyiJorgensenKarpelSanadhya2025, BezuglyiKarpelKwiatkowskiWata2024,BezuglyiJorgensenKarpelKwiatkowski2025, BezuglyiDudkoKarpel2024, BezuglyiJorgensenKarpelSanadhyaRaszeja2026}). 
Though Bratteli diagrams provide a convenient framework for studying dynamical systems in Cantor and Borel dynamics, there is no unique way to associate a diagram to a transformation. 
In general, it remains an open problem to determine whether given two Bratteli diagrams model isomorphic dynamical systems (see, e.g., \cite{DownarowiczKarpel_2019}). 
%\tcr{more references?} 
In this paper, we consider the natural notion of isomorphism for Bratteli diagrams, which can be found, for instance, in  \cite[Section 2]{HermanPutnamSkau1992} or \cite[Subsection 6.3.1]{Durand2010} for standard Bratteli diagrams, and in \cite{BezuglyiJorgensenKarpelSanadhya2025} for generalized ones. 
The diagrams isomorphic with respect to this natural isomorphism have Borel isomorphic tail equivalence relations. 
%\tcr{add this later too, in Section 4? Even homeomorphic relations if it makes sense? It was some open question about continuous isomorphisms?} 

%We classify (generalized) Bratteli diagrams with respect to isomorphism and
We study the invariants of isomorphism of generalized Bratteli diagrams.
%study its invariants.
%surprisingly different shape, study invariants... irreducibility ... 
%difference from usual graph theory
%difference from standard BD
%This paper is devoted to the classification of generalized Bratteli diagrams with respect to isomorphism.
%The natural notion of isomorphism between standard Bratteli diagrams is mentioned, for instance, in \cite[Section 2]{HermanPutnamSkau1992} or \cite[Subsection 6.3.1]{Durand2010}. 
%Two Bratteli diagrams $B = (V,E)$ and $B' = (V',E')$ are isomorphic if there are  bijections $(g_n \colon V_n  \rightarrow V'_n)_{n}$, and $(h_n \colon E_n \rightarrow E'_n)_n$ which intertwine the corresponding source and range maps. 
%This isomorphism is given by a pair of bijections between the sets of vertices and between the sets of edges of two Bratteli diagrams, preserving the gradings and intertwining the corresponding source and range maps. 
%The same notion of isomorphism for generalized Bratteli diagrams is briefly discussed in \cite{BezuglyiJorgensenKarpelSanadhya2025}. 
%In this paper, we systematically study this notion, focusing on the case of generalized Bratteli diagrams. 
%We discuss which properties of Bratteli diagrams are preserved under isomorphism. 
In particular, we focus on the notion of irreducibility. For a Bratteli diagram, the set of vertices is partitioned into subsets $V = \bigsqcup_{n = 0}^{\infty} V_n$ called levels. The sets of vertices of each level of a generalized Bratteli diagram is countably infinite, each set $V_n$ can be identified with the same countably infinite set, e.g., the set of natural numbers or the set integers. 
The diagram is called irreducible if 
for any vertices $i, j \in V_0$ and any $n \in \N_0$ there exist 
$m 
> n$ such that there is a finite path between $i \in V_n$ and $j \in V_m$.
%We introduce the notion of completely irreducible generalized Bratteli diagrams.
%\tcr{Is it Ok to suddenly start talking about details here, not explaining anything earlier about vertices, levels, etc? Maybe I should add a bit more explanation earlier? And that for standard BD the notion of irreducibility only makes sense for finite rank? Am I correct?}
In general, irreducibility is not preserved under isomorphism, but we show that there are generalized Bratteli diagrams
which are isomorphic only to irreducible generalized Bratteli diagrams. We call such diagrams completely irreducible.
%always irreducible independently of the enumeration of their vertices on each level. 
Complete irreducibility implies minimality of the tail equivalence relation, thus, completely irreducible generalized Bratteli diagrams can be viewed as analogues of simple standard Bratteli diagrams.
%\tcr{write for fin rank}
%\tcr{Maybe here I should add more that simple BD are not possible in the generalized case?}
We give necessary conditions and sufficient conditions  for a generalized Bratteli diagram to be completely irreducible. 
We prove that every  generalized Bratteli diagram is isomorphic to an irreducible generalized Bratteli diagram. 
We show that irreducibility and connectedness as an undirected graph are independent properties for generalized Bratteli diagrams. We prove that topological transitivity of the tail equivalence relation implies connectedness of the Bratteli diagram, though the converse statement is not true.
The results are illustrated by numerous examples. We note that a similar notion of irreducibility can also be introduced for finite rank standard Bratteli diagrams. The notion of irreducibility is also known for Markov chains, although the tail equivalence relation is not considered there (see e.g. \cite{Seneta2006}).

%\tcr{write about tail equiv relation and incid matrices}
%\tcr{we introduce another isom notion}

%\tcr{highlight differences between standard and gbd}


%Irreducibility, isomorphism, top prop of tail equivalence relation, connectedness.

%\tcr{natural q, classification problems, isom ref HPS etc} 
%The notion of isomorphism of generalized Bratteli diagrams was introduced in \cite{BezuglyiJorgensenKarpelSanadhya2025}. 
%\tcr{Durand survey see + general graded graph isom}



%We connect the notion of irreducibility of a generalized Bratteli diagram with the properties of the corresponding tail equivalence relation.

%\tcr{what's known for standard stationary BD...}

%Sometimes it is possible to introduce an order on a standard Bratteli diagram $B$ which generates a homeomorphism of the path space $X_B$ called a Vershik map (see e.g. \cite{BezuglyiKwiatkowskiYassawi2014}, \cite{BezuglyiYassawi2017}).

%Simple - always prop order and homeo...

%Topological models of Borel automorphisms (as homeomorphisms / continuous maps of locally compact Cantor spaces / Polish spaces).

%wandering dynamics, smooth Borel eq rel on BD?

%\tcr{continuous isomorphism vs not contin?}

%\tcr{what's known for standard BD / graphs??}

%\tcr{Add about usage of BD from motivation slides}

%\tcb{It was proved in \cite{BezuglyiDooleyKwiatkowski2006} that every aperiodic hyperfinite countable Borel equivalence relation can be represented as a tail equivalence relation on a generalized Bratteli diagram. }

%\tcr{add literature about sds on infinite alphabets, Ferenczi}





%\tcr{add shortly abt simple and irred bd}

%\tcr{maybe move to section 3}

%for which the tail equivalence relation is minimal. \tcr{is it in both directions?}


%Generalized Bratteli diagrams of bounded size have a locally compact path space and hence can be considered to be the class of generalized Bratteli diagrams which is most reminiscent of standard Bratteli diagrams (see \cite[Section 3]{BezuglyiJorgensenKarpelSanadhya2025}). Generalized Bratteli diagrams of bounded size also provide models for substitution dynamical systems on countably infinite alphabets (see \cite{BezuglyiJorgensenSanadhya2024}).

%\tcr{look at BDK 06: loc comp gbd, what about v0}

%\tcb{An irreducible countable Markov chain is a stochastic process with a countably infinite or finite state space where every state can reach every other state in a finite number of steps.}

%It was proved in \cite[Theorem 5.1]{BezuglyiJorgensenKarpelSanadhya2025} that a generalized stationary Bratteli diagram with an irreducible aperiodic incidence matrix has a topologically transitive tail equivalence relation. \tcr{see remark 5.2 there! irreducibility vs existence of vertical paths}


%\tcr{add abt V map, CBER from first slides}
%\tcr{question naturally arises, first step...}

%The structure of the diagram ... It was proved in \cite[Theorem 5.1]{BezuglyiJorgensenKarpelSanadhya2025} that the tail equivalence relation on the path space of a stationary generalized Bratteli diagram with an irreducible and aperiodic incidence matrix is topologically transitive. It was also shown in  \cite[Theorem 5.4]{BezuglyiJorgensenKarpelSanadhya2025} that the irreducibility of the diagram does not imply that the tail equivalence relation is minimal. \tcr{add remark from that paper also later}

%\tcb{ In the studies of invariant measures, one is interested in the incidence matrices of the Bratteli diagram. Thus, in this paper, we will be interested in the way the incidence matrices of a generalized Bratteli diagram can change under the isomorphism and how properties of the tail equivalence relation influences the isomorphism classes of generalized Bratteli diagrams.}

%\tcb{We will refer to the bijection $g_n$ as ``enumeration'' of vertices on level $n$. From Definition~\ref{def_isom BD} it follows that for $v \in V_n$ and $w \in E_m$, we have $|E(v,w)| = |E(g_n(v), g_m(w))|$. Since we will be interested in incidence matrices and properties of the tail equivalence relation, ... edge isom doesn't matter}

%\tcr{paper with Jan beginning enumeration a(n)}
%\tcr{def tail equiv relation}

The outline of the paper is as follows. In Section~\ref{sec:prelim}, we provide the notation and basic definitions that we  use throughout the paper. In Section \ref{sec:examples}, we study in detail several classes of generalized Bratteli diagrams which demonstrate various behaviours in terms of irreducibility and topological properties of the corresponding tail equivalence relation, such as minimality and topological transitivity. These diagrams also
serve later as examples for the results from Section \ref{sec:mainres}. Section \ref{sec:mainres} contains the main results of the paper. We introduce the notions of completely irreducible and relatively irreducible generalized Bratteli diagrams. We show that every generalized Bratteli diagram is either completely or relatively irreducible. We provide necessary conditions and sufficient conditions for a Bratteli diagram to be completely irreducible. %We provide examples showing that irreducibility and connectedness are independent properties for generalized Bratteli diagrams. 
We study the relations between connectedness, irreducibility, and topological properties of the tail equivalence relation, and of the path space of a generalized Bratteli diagram.
Our main results are contained in Theorems \ref{thm:classII}, \ref{Thm:isomredus}, \ref{thm:Class1_minimalR}.



%We provide some necessary and sufficient conditions for a generalized Bratteli diagram to belong to these classes. 


%\tcr{add connections to other papers on substitutions}

%\tcr{BJKS Irreducibility + vertical paths imply conditions on the tail equiv relation}

%Class I is analogue of notion of simplicity for standard BD: add examples 0-1 on rank 2 bd 

\section{Preliminaries}\label{sec:prelim}

This section contains the main definitions concerning generalized Bratteli diagrams.  
%The class of generalized Bratteli diagrams was defined in \cite{BezuglyiDooleyKwiatkowski_2006} and then studied in recent publications \cite{BezuglyiJorgensen2022}, \cite{BezuglyiJorgensenKarpelSanadhya2025}, \cite{Bezuglyi_Jorgensen_Sanadhya_2023}, \cite{BezuglyiKarpelKwiatkowski2024}, see also \cite{BezuglyiKarpelKwiatkowskiWata2024}, \cite{BezuglyiJorgensenKarpelKwiatkowski2025}, \cite{BezuglyiDudkoKarpel2024}.
In particular, we discuss the notions of isomorphism and irreducibility for generalized Bratteli diagrams. For more details on generalized Bratteli diagrams see, e.g., \cite{BezuglyiJorgensenKarpelSanadhya2025} or \cite{BezuglyiKarpelKwiatkowskiWata2024}.

We use the standard notation $\mathbb N, \Z, \R,$  
$\mathbb{N}_0 = \N \cup \{0\}$ for the sets of numbers, 
$|\cdot |$ denotes the cardinality of a set.

%\subsection{Main definitions}
%\subsection{Standard and generalized Bratteli diagrams}


\begin{definition}\label{Def:generalized_BD} A 
\textit{(standard or generalized) Bratteli diagram} is a graded graph 
$B = (V, E)$ such that the 
vertex set $V$ and the edge set $E$ can be partitioned $V = \bigsqcup_{n=0}^\infty  V_n$ and $E = 
\bigsqcup_{n=0}^\infty  E_n$ so that the following 
properties hold: 

\begin{enumerate}[label=\upshape(\roman*), leftmargin=*, widest=iii]

\item for every $n \in \mathbb{N}_0$, the number of vertices at each level 
$V_n$ is countable, it is finite for standard Bratteli diagrams and countably infinite for generalized Bratteli diagrams. Similarly, the set $E_n$
of all edges between $V_n$ and $V_{n+1}$ is countable and it is finite for standard Bratteli diagrams and countably infinite for generalized Bratteli diagrams,

\item for every edge $e\in E$, we define the 
range and 
source maps $r$ and $s$ such that $r(E_n) = V_{n+1}$ and 
$s(E_n) = V_{n}$ for $n \in \N_0$,

\item for every vertex $v \in V \setminus V_0$, we 
have $|r^{-1}(v)| < \infty$ (the condition holds automatically for standard Bratteli diagrams).
\end{enumerate}
\end{definition} 



\begin{remark}
For all $n$, the edges of $E_n$ always have sources in $V_{n}$ and ranges in $V_{n+1}$, hence in the pictures all the edges should be directed downwards. For the better readability of figures, we draw all Bratteli diagrams as unordered graphs.
\end{remark}

\begin{remark}
    For standard Bratteli diagrams, there is a convention which comes from the paper of O. Bratteli \cite{Bratteli1972} to assume that $V_0 = \{v_0\}$ is a single point. Since this paper mostly focuses on the case of generalized Bratteli diagrams, for the sake of the uniformity of notation, we drop this assumption and allow $V_0$ to be arbitrary finite set in the case of standard Bratteli diagrams. 
\end{remark}

Let $B$ be a (standard or generalized) Bratteli diagram.
The set $V_n$ is called the \textit{$n$th level} of the diagram $B$. For generalized Bratteli diagrams, we usually identify each $V_n$ either with $\N$ (or $\N_0$) or with $\Z$. If the vertices of each level of a generalized Bratteli diagram are enumerated by $\N$ or $\N_0$, we call such diagram \textit{one-sided}, if the the vertices are enumerated by $\Z$, the diagram is called \textit{two-sided}. 


% A sequence of edges $(e_m, \ldots, e_n)$, where $m < n$ and $e_i \in E_i$ for $i = m, \ldots, n$ form a finite path if $r(e_{j-1}) = s(e_j)$ for all $j = m+1, \ldots, n$
To define the \textit{path space} of a generalized Bratteli diagram
$B$, consider 
a (finite or infinite) sequence of edges
 $ x = (e_i: e_i\in E_i)$ (it is called a \textit{path}) such 
that $s(e_i)=r(e_{i-1})$. Denote the set of all infinite
paths $x$ starting at a vertex in $V_0$ by $X_B$ and call 
it the \textit{path space} of the diagram $B$. For a finite path
$\ol e = (e_0, ... , e_n)$, let 
$s(\ol e) = s(e_0)$ and $r(\ol e) = r(e_n)$. The set
$$
    [\ol e] := \{x = (x_i) \in X_B : x_0 = e_0, ..., x_n = e_n\}, 
$$ 
is called the \textit{cylinder set} associated with $\ol e$.



The \textit{topology} on the path space $X_B$ is generated by
cylinder sets that are clopen in this topology.
This topology coincides with the topology defined by the 
following metric on $X_B$: for $x = (x_i), \, y = (y_i)$, set 
$$
\mathrm{dist}(x, y) = \frac{1}{2^N},\ \ \ N = \min\{i \in \N_0 : 
x_i \neq y_i\}.
$$
The path space $X_B$ is a zero-dimensional Polish space and, therefore, a standard Borel space. In general, $X_B$ is not locally compact. 
%As usual, we will consider the Bratteli diagrams whose path spaces have no isolated points.


A matrix $A = (a_{ij})$ is called \textit{infinite} 
(or countably infinite) if its rows and columns
are indexed by the same countably infinite set.
For vertices $v \in V_m$ and $w \in V_{n}$, 
denote 
by $E(v, w)$ the set of all finite paths between $v$ and $w$. 
Set $f^{(n)}_{vw} = |E(v, w)|$ for every $w \in V_n$ and $v \in 
V_{n+1}$. In other words, $f^{(n)}_{vw}$ is a number of edges between the vertices $w \in V_n$ and $v \in V_{n+1}$. In such a way,  we associate with the generalized 
Bratteli diagram $B = (V,E)$ 
a sequence of non-negative countably infinite matrices $(F_n)$, 
$n \in \N_0$, 
called the \textit{incidence  matrices}: 
\begin{equation}\label{Notation:f^n}
    F_n = (f^{(n)}_{vw} : v \in V_{n+1}, w\in V_n),\ \   
    f^{(n)}_{vw}  \in \N_0.
\end{equation}

We will write $B = B(F_n)$. The structure of a (standard or generalized) generalized Bratteli diagram $B=(V, E)$ is 
completely determined by the sequence of its incidence matrices 
$(F_n), 
\, n \in \N_0$.  For each $n 
\in \N_0$, the matrix $F_n$ has at most finitely many non-zero 
entries in each row, and none of its rows or columns are entirely 
zero. For a generalized Bratteli diagram $B$, a column of $F_n$ may have a finite or infinite number of 
non-zero entries. 
%The path space $X_B$ is locally compact if and only if every column of $F_n$ has finitely many non-zero entries for all $n$ \tcr{not true?}. 
If  $F_n = F$ for every $n \in 
\N_0$, then the diagram $B$ is called \textit{stationary.} We 
will write $B = B(F)$ in this case.



\begin{definition}\label{Def:Tail_equiv_relation}
Two paths $x= (x_i)$ and $y=(y_i)$ in $X_B$ are called 
\textit{tail equivalent} if there exists an $n \in \mathbb{N}_0$ 
such that $x_i = y_i$ for all $i \geq n$. This notion defines a \textit{countable Borel equivalence relation} $\mathcal R$ in the path space $X_B$, which is called the \textit{tail equivalence relation}.
\end{definition}

Denote by $\mathcal{O}(x)$ the orbit of a point $x \in X_B$ under the tail equivalence relation:
$$
\mathcal{O}(x) = \{y \in X_B : (x,y) \in \mathcal{R}\}.
$$
The tail equivalence relation is called \textit{minimal} if the orbit of every point is dense, and it is called \textit{topologically transitive} if there exists a point with a dense orbit (see e.g. \cite{KatokHasselblatt1995} or \cite{Bruin2022}). %\tcr{refs here?}

The following notion of isomorphism can be found, for instance, in \cite[Section 2]{HermanPutnamSkau1992} and \cite[Subsection 6.3.1]{Durand2010} for standard Bratteli diagrams, and in \cite[Definition 2.8]{BezuglyiJorgensenKarpelSanadhya2025} for generalized Bratteli diagrams.

\begin{definition} \label{def_isom BD}
Two (standard or generalized) Bratteli diagrams $B = (V, E)$ and $B' = 
(V', E')$ 
are called \textit{isomorphic} if there exist two sequence of
bijections $(g_n : V_n \rightarrow V_n')_{n \in \N_0}$ and 
$(h_n : E_n \rightarrow E_n')_{n \in \N_0}$ such that 
for every $n \in \N_0$, we have $g_n(V_n) = V_n'$ and $h_n(E_n) = 
E_n'$, and $s' \circ h_n = g_n \circ s$, $r' \circ h_n = 
g_n \circ r$,  where $s'$ and $r'$ are source and range maps in $B'$.
\end{definition}

In most of the studies concerning standard and generalized Bratteli diagrams, one is interested in the structure of the diagram, which is completely determined by the sequence of incidence matrices. 
%For instance, it is enough to give the sequence of incidence matrices to study the tail equivalence relation or tail invariant measures (ref)
Thus, in this paper, we focus on the following definition rather than on the definition of isomorphism:

%\tcr{don't care about edges: glue isomorphisms, consider equivalence classes...}

\begin{definition}
Let $B = (V,E)$ and $B' = (V',E')$ be two (standard or generalized) Bratteli diagrams with the corresponding sequences of incidence matrices $(F_n)$ and $(F_n')$. Assume that for every $n \in \mathbb{N}_0$ there exists a sequence of bijections $g_n \colon V_n \rightarrow V'_n$ such that $g_n(V_n) = V'_n$ and $f^{(n)}_{vw} = {f'}^{(n)}_{g_{n+1}(v),g_{n}(w)}$. Then we say that the diagrams $B$ and $B'$ are \textit{isomophic with respect to the sequence of vertex bijections $(g_n)$} and denote $B \simeq B'$. 
\end{definition}


If the incidence matrices of $B$ and $B'$ are $0-1$ matrices, then the usual definition of isomorphism and the definition of isomorphism with respect to the sequence of vertex bijections coincide. If, for some $n \in \N_0$, the vertices $v \in V_{n+1}$ and $w \in V_n$ in the diagram $B$ are connected by multiple edges, then we are not concerned with how the set of edges $E(v, w)$ is mapped into the set of edges $E(g_{n+1}(v), g_n(w))$. In other words, if $B \simeq B'$, there may exist multiple isomorphisms mapping $B$ to $B'$. Clearly, $\simeq$ is an equivalence relation.

%Since in this paper we will only consider the notion of isomophism with respect to the sequence of vertex bijections, to simplify the notation we will call Bratteli diagrams $B$ and $B'$ isomo


%In other words, if there are  we are not interested in the particular way the edges $E_n$ are mapped into $E'_n$

%Note that if $B \simeq B'$ then they are isomorphic according to Definition \ref{def_isom BD}. Indeed, as $h_n \colon E_n \rightarrow E_n'$

Below we formulate the fact that the diagrams $B$ and $B'$ are isomorphic in terms of their incidence matrices.
By a \textit{permutation matrix} we will mean a matrix that has exactly one entry of $1$ in each row and each column with all other entries equal to $0$. If $P$ is a permutation matrix, then it is easy to see that $PP^T = P^TP = I$, where $I = (I_{kl})$ is a (finite or infinite) square matrix with $I_{kl} = 1$ for $k = l$ and $I_{kl} = 0$ for $k \neq l$. Clearly, if the diagram $B$ is isomophic to the diagram $B'$ with respect to the sequence of vertex bijections then the following result holds.
% An infinite square matrix is a linear operator represented by an array with an infinite number of rows and columns, typically acting on infinite-dimensional vector spaces. These matrices, often indexed by integers, are used in quantum mechanics (Hilbert spaces) and differential equations, where they can represent operators like derivatives (AI review)

\begin{prop}\label{thm:isomBD}
    Let $B= B(F_n)$ and $B' = B'(F'_n)$ be (standard or generalized) Bratteli diagrams. Then $B \simeq B'$ if and only if there exist permutation matrices $P_n$ such that for all $n \in \N_0$:
    $$
    F'_n = P_n F_n P_{n+1}^T.
    $$
\end{prop}


To illustrate Definition \ref{def_isom BD} and Proposition \ref{thm:isomBD}, we consider the following example. This example can also be found in \cite[Example 2.9]{BezuglyiJorgensenKarpelSanadhya2025}. Here, we provide an explicit description of permutation matrices which are used for the isomorphism.



%The following example illustrates the notion of isomorphism of generalized Bratteli diagrams and can be also found in \cite[Example 2.9]{BezuglyiJorgensenKarpelSanadhya2025}. Here we also add the explicit description of permutation matrices which are used for the isomorphism.



\begin{example}[Isomorphic generalized Bratteli 
diagrams, Figure~\ref{Fig:Isom_BD}]\label{Ex:isom_bd}
Let $B(F)$ be a stationary generalized Bratteli diagram such that 
every level of $B$ is identified with $\Z$, and the incidence 
matrix $F = (f_{ij})_{i,j \in \Z}$ has entries
$$
f_{ij} = 
\left\{
\begin{aligned}
& 2, \mbox{ for } i = j,\\
& 1, \mbox{ for } |i - j| = 1,\\
& 0, \mbox{ otherwise. }
\end{aligned}
\right.
$$ 
Let $B'(F')$ be a stationary generalized Bratteli 
diagram such that every level of $B'$ is identified with $\N_0$, 
and its incidence matrix $F' = (f'_{ij})_{i,j \in \N_{0}}$ has 
entries $f'_{00} = f'_{11} = 2$, $f'_{01} = 1, f'_{02} = 1, f'_{10} = 1, f'_{13} = 1, f'_{20} = 1, f'_{31} = 1$,
and for $i,j \notin \{0,1\}$: 
$$
f'_{ij} = 
\left\{
\begin{aligned}
& 2, \mbox{ for } i = j,\\
& 1, \mbox{ for } |i - j| = 2,\\
& 0, \mbox{ otherwise. }
\end{aligned}
\right.
$$ 
The diagrams $B$ and $B'$ (elaborated in 
Figure~\ref{Fig:Isom_BD}) are isomorphic. Indeed, the bijections 
$(g_n : V_n \rightarrow V_n')_{n \in \N_0}$ and $(h_n : E_n 
\rightarrow E_n')_{n \in \N_0}$ that give the isomorphism are 
defined as follows: since the diagrams are stationary, we set,
for every $n \in \N_0$, $g_n = g : \mathbb{Z} \rightarrow 
\mathbb{N}_0$ 
$$
g (n) = 
\left\{
\begin{aligned}
&2n, \mbox{ if } n \geq 0,\\
&- 2n - 1, \mbox{ if } n < 0.
\end{aligned}
\right.
$$ 
The bijection $h_n : E_n \rightarrow E_n'$ is defined as follows: 
the two vertical edges in the diagram $B(F)$ with range $i \in 
V_n$ are mapped to the two vertical edges in the diagram $B'(F')$ 
with range $g(i) \in V_n$. For $i > 0$, the 
non-vertical edge with range $i$ coming from left (respectively 
from the right) is mapped to the non-vertical edge 
with range 
$g(i) \in V'_n$ coming from the left (respectively from the 
right). For $i < 0$, the 
non-vertical edge with range $i$ coming from left (respectively 
from the right) is mapped to the non-vertical edge 
with range 
$g(i) \in V'_n$ coming from the right (respectively from the 
left). For the non-vertical edges with range $0 \in V_n$, 
the mapping  can be seen from the 
Figure~\ref{Fig:Isom_BD}. It is easy to see that with the above 
bijections the two diagrams are isomorphic. 



\begin{figure}
\unitlength=1cm
\begin{graph}(11,4)
% \graphnodesize{0.2}
% \roundnode{V0}(3,6)
%  %\nodetext{V0}(-1,0){$V_0$}
 \roundnode{V11}(2,3)
 %\nodetext{V11}(-0.7,0){$w_1^{(0)}$}
  %\nodetext{V12}(0.7,0){$w_2^{(0)}$}
 \roundnode{V12}(4,3)
 \roundnode{V13}(6,3)
 \roundnode{V14}(8,3)
  \roundnode{V15}(10,3)
  % The second level vertices
 \roundnode{V21}(2,1)
 \roundnode{V22}(4,1)
 \roundnode{V23}(6,1)
  \roundnode{V24}(8,1)
    \roundnode{V25}(10,1)
  %\nodetext{V21}(-0.7,0){$w_1^{(1)}$}
 % \nodetext{V22}(0.7,0){$w_2^{(1)}$}
  %
 %
 % EDGES
 \graphlinewidth{0.025}
% % First level
%  \edge{V0}{V11}
%  \edge{V0}{V12}

 % Second level
 \bow{V21}{V11}{0.09}
  \bow{V21}{V11}{-0.09}
    \edge{V22}{V11}
    \edge{V21}{V12}
     
 \bow{V22}{V12}{0.09}
 \bow{V22}{V12}{-0.09}

 
 \bow{V23}{V13}{0.09}
  \bow{V23}{V13}{-0.09}

 \edge{V23}{V12}
    \edge{V22}{V13}
    
     \bow{V24}{V14}{0.09}
  \bow{V24}{V14}{-0.09}
  
   \edge{V24}{V13}
    \edge{V23}{V14}
    
         \bow{V25}{V15}{0.09}
  \bow{V25}{V15}{-0.09}
  
   \edge{V25}{V14}
    \edge{V24}{V15}

 
 \freetext(10.9,3){$\ldots$}
  \freetext(10.9,1){$\ldots$}
   \freetext(1,3){$\ldots$}
  \freetext(1,1){$\ldots$}
  \freetext(2,0.5){$\vdots$}
  \freetext(4,0.5){$\vdots$}
    \freetext(6,0.5){$\vdots$}
      \freetext(8,0.5){$\vdots$}  
        \freetext(10,0.5){$\vdots$} 
    
\end{graph}

\begin{graph}(11,4)
% \graphnodesize{0.2}
% \roundnode{V0}(3,6)
%  %\nodetext{V0}(-1,0){$V_0$}
 \roundnode{V11}(2,3)
 %\nodetext{V11}(-0.7,0){$w_1^{(0)}$}
  %\nodetext{V12}(0.7,0){$w_2^{(0)}$}
 \roundnode{V12}(4,3)
 \roundnode{V13}(6,3)
 \roundnode{V14}(8,3)
  \roundnode{V15}(10,3)
  % The second level vertices
 \roundnode{V21}(2,1)
 \roundnode{V22}(4,1)
 \roundnode{V23}(6,1)
  \roundnode{V24}(8,1)
    \roundnode{V25}(10,1)
  %\nodetext{V21}(-0.7,0){$w_1^{(1)}$}
 % \nodetext{V22}(0.7,0){$w_2^{(1)}$}
  %
 %
 % EDGES
 \graphlinewidth{0.025}
% % First level
%  \edge{V0}{V11}
%  \edge{V0}{V12}

 % Second level
 \bow{V21}{V11}{0.09}
  \bow{V21}{V11}{-0.09}
    \edge{V22}{V11}
    \edge{V21}{V12}
     \edge{V21}{V13}
     
 \bow{V22}{V12}{0.09}
 \bow{V22}{V12}{-0.09}

 
 \bow{V23}{V13}{0.09}
  \bow{V23}{V13}{-0.09}

 \edge{V23}{V11}
    \edge{V22}{V14}
    
     \bow{V24}{V14}{0.09}
  \bow{V24}{V14}{-0.09}
  
   \edge{V24}{V12}
    \edge{V23}{V15}
    
         \bow{V25}{V15}{0.09}
  \bow{V25}{V15}{-0.09}
  
   \edge{V25}{V13}
    %\edge{V24}{V15}

 
 \freetext(10.9,3){$\ldots$}
  \freetext(10.9,1){$\ldots$}
  \freetext(2,0.5){$\vdots$}
  \freetext(4,0.5){$\vdots$}
    \freetext(6,0.5){$\vdots$}
      \freetext(8,0.5){$\vdots$}  
        \freetext(10,0.5){$\vdots$} 
    
\end{graph}
\caption{Isomorphic generalized Bratteli diagrams $B$ and 
$B'$.}\label{Fig:Isom_BD}
\end{figure}

We can formulate the fact that the diagrams $B$ and $B'$ are isomorphic in terms of their incidence matrices. Let $P$ be a $\mathbb{N}_0 \times \mathbb{Z}$ matrix such that
$$
p_{ij} = 
\left\{
\begin{aligned}
& 1, \mbox{ for } i = g(j),\\
& 0, \mbox{ otherwise. }
\end{aligned}
\right.
$$ 
It is straightforward to verify that $F' = PFP^T$.
\end{example}
%doublecheck


%\tcr{Mention V n either naturals or integers, OP: consider other enumerations}

%\tcb{We will refer to the bijection $g_n$ as ``enumeration'' of vertices on level $n$. From Definition~\ref{def_isom BD} it follows that for $v \in V_n$ and $w \in E_m$, we have $|E(v,w)| = |E(g_n(v), g_m(w))|$. Since we will be interested in incidence matrices and properties of the tail equivalence relation, ... edge isom doesn't matter}

%\tcr{paper with Jan beginning enumeration a(n)}
%\tcr{def tail equiv relation}

%\tcr{isomorphism - enumeration of vertices, only important the number of edges (for the study of measures etc incidence matr important)}

\begin{definition}\label{Def:irreducible_GBD} 
Let $B = B(F_n)$ be
a generalized Bratteli diagram. Assume that all levels $V_i$ of $B$ are identified with a set $V_0$ (e.g. $V_0 = 
\mathbb{N}$ or $\mathbb{Z}$). Then $B$ is called  \textit{irreducible} if 
for any vertices $i, j \in V_0$ and any $n \in \N_0$ there exist 
$m 
> n$ such that there is a finite path between $i \in V_n$ and $j \in V_m$. In 
other words, the $(j, i)$-entry of the matrix $F_{m-1} \cdots 
F_n$ 
is non-zero. Otherwise, the diagram is called \textit{reducible}.
\end{definition}


\begin{remark}
Similarly to the case of finite matrices, a countably infinite matrix $A = (a_{ij})$ is called 
\textit{irreducible} if for every pair of indices $i,j$ there 
is $n > 0$ such that $a^{(n)}_{ij} > 0$, where $a^{(n)}_{ij}$ are elements of the matrix $A^n$. 
A stationary generalized Bratteli diagram is
irreducible if and only if the corresponding incidence matrix is 
irreducible.
Denote $$
p(i) = \gcd\{n:a^{(n)}_{ii} > 0\}.
$$
Then $p(i)$ is called the \textit{period of index} $i$. For an 
irreducible matrix $A$, the periods of all indices are the same 
and called the \textit{period of $A$}. An irreducible matrix with 
period one is called \textit{aperiodic} (see \cite{Kitchens1998} for 
more definitions and results about infinite non-negative matrices).
\end{remark}
%Recall for every level $n \in \mathbb{N}_0$, the set of vertices $V_n$ of a generalized Bratteli diagram is identified with the same set (usually $\mathbb{N}$ or $\mathbb{Z}$). 

\begin{remark}
    We can introduce similar notion of irreducibility for standard Bratteli diagrams, which have the same number of vertices on each level. Such diagrams are diagrams of finite rank (see \cite{BezuglyiKwiatkowskiMedynetsSolomyak2013}). 
\end{remark}

Recall that we identify the vertices of each level $V_n$ of a generalized Bratteli diagram $B$ with the same countably infinite set, usually $\N$ or $\Z$. 
After such identification,  we will call the edges $e_n$ for which $s(e_n) = i \in V_n$ and $r(e_n) = i \in V_{n+1}$ \textit{vertical} and the edges with $s(e_n) = i \in V_n$ and $r(e_n) = j \in V_{n+1}$ with $i \neq j$ \textit{slanted}. %\tcr{slanted or slanting?}
We call an infinite path \textit{vertical} if all its edges are vertical and \textit{slanted} if all its edges are slanted. It was noticed in \cite[Remark 2.10]{BezuglyiJorgensenKarpelSanadhya2025} that, in general, the property of irreducibility is not preserved under an isomorphism of generalized Bratteli diagrams. For the readers convenience, we repeat the remark below.


\begin{remark}\label{Rem:isomvsirred}\cite{BezuglyiJorgensenKarpelSanadhya2025}
Note that, in general, isomorphism of generalized Bratteli diagrams does not preserve irreducibility. For instance, stationary diagram $B$ from Example~\ref{Ex:isom_bd} is irreducible since starting from every vertex we eventually reach any other vertex for some sufficiently low level. It is easy to see that diagram $B'$ is also irreducible.
    Define another stationary generalized Bratteli diagram 
$B''=(V'', E'')$ isomorphic to $B$ with every level identified with $\Z$. Consider the following sequence of  bijections $g_n \colon V_n \rightarrow V_n''$ between vertices of $B$ and $B''$ on level $n$:
    $$
    g_n(i) = i + n.
    $$
    Thus, the bijections $g_n$ do not change the level $V_0$, shift all the vertices of level $V_1$ by $1$ to the right and all vertices of level $V_n$ by $n$ to the right. 
    For all $n$, the corresponding sequence of bijections $h_n \colon E_n \rightarrow E''_n$ between edges will map all edges outgoing from vertex $i \in V_n$ to the edges with the source $i + n \in V'_n$ and the range 
   in the set $\{i+n, i+n+1, i+n+2\}.$ 
In particular,
the edge with source $i \in V_n$ and range $i - 1 \in V_{n+1}$ will be mapped to the vertical edge between $i+n \in V''_n$ and $i+n \in V''_{n+1}$. Thus, for every $n \in \N_0$, it will not be possible to reach from vertex $i \in V''_n$ any vertex $j < i$ for any level $m > n$. The incidence matrix $F''$ of $B''$ is lower triangular (here we indicate the main diagonal of $F''$ with bold font):
    $$
    F'' = 
    \begin{pmatrix}
       \ddots & \vdots & \vdots & \vdots & \vdots & \udots\\
       \ldots & \textbf{1} & 0 & 0 & 0 &\ldots\\
       \ldots & {2} & \textbf{1} & 0 & 0 & \ldots\\   
       \ldots & {1} & 2 & \textbf{1} & 0 & \ldots\\ 
       \ldots & 0 & 1 & 2 & \textbf{1} & \ldots\\     
       \udots & \vdots & \vdots & \vdots & \vdots & \ddots
    \end{pmatrix}.
    $$
\end{remark}

%\tcr{where to put it? Write about abuse of notation}
%\medskip


\section{Classes of generalized Bratteli diagrams}\label{sec:examples}

In this section, we examine in detail several classes of generalized Bratteli diagrams that exhibit a range of behaviours with respect to irreducibility and the topological properties of their tail equivalence relations.
These diagrams also serve as examples illustrating the results presented in Section \ref{sec:mainres}.


\subsection{Generalized Bratteli diagrams of bounded size}\label{subsec:bdd}

Generalized Bratteli diagrams of bounded size  
are used to 
model substitution dynamical systems on infinite alphabets 
(for more information on substitutions on infinite alphabets see, e.g., \cite{Ferenczi_2006}, \cite{DomingosFerencziMessaoudiValle2024}, \cite{Frettloh_Garber_Manibo_2022}, and for the corresponding Bratteli diagram construction see \cite{BezuglyiJorgensenSanadhya2024}). %These diagrams are characterized by the fact that all incidence matrices are banded. 
Such diagrams have a locally compact path space and can be considered as an intermediate step between 
the standard and generalized Bratteli diagrams. 

%when they are irreducible, top transitive tail equivalence relation

%\tcr{here or in intro, add references from GGD}



\begin{definition}\label{Def:BD_bdd_size} A generalized Bratteli 
diagram $B(F_n)$ is called of \textit{bounded size} if for all $n \geq 0$ the set $V_n$ is identified with $\mathbb{Z}$ and there exists 
a sequence of pairs of natural numbers $(t_n, L_n)_{n \in \N_0}$ 
such that, for all $n \in \mathbb{N}_0$ and all $v \in V_{n+1}$, we have
\begin{equation}\label{eq: Bndd size}
s(r^{-1}(v)) \in \{v - t_n, \ldots, v + t_n\} \quad \mbox{and} 
\quad \sum_{w \in V_{n}} f^{(n)}_{vw} = \sum_{w \in V_{n}} |E(w,v)| 
\leq L_n.
\end{equation} 
%If the sequence $(t_n, L_n)_{n \in \N_0}$ is constant, i.e. $t_n = t$ and $L_n = L$ for all $n \in \N_0$, then we say that the diagram $B(F_n)$ is of \textit{uniformly bounded size}.
\end{definition} 

\begin{remark}
To simplify the notation and calculations, we assume that for each $n \in \N_0$, the natural numbers 
$(t_n, L_n)$ are chosen to be the minimal possible. Also, for every 
$n \in \mathbb{N}_0$, it is assumed that $E(v - t_n, v)$ and 
$E(v + t_n, v)$ are nonempty for all $v \in V_{n+1}$. Throughout this paper, we will not use the bound given by $L_n$ and will be using only the sequence of parameters $(t_n)$.
\end{remark}

\begin{remark}
    In this paper, we will only use the parameters $(t_n)$ for the generalized Bratteli diagrams of bounded size. In general, whether there is a single edge or multiple edges between a given pair of vertices doesn't affect irreducibility or connectedness of a Bratteli diagram, or whether the tail equivalence relation is topologically transitive or minimal. %\tcr{move earlier?}
\end{remark}

Let $B$ be a generalized Brattlei diagram of bounded size. Since every vertex of $B$ has finitely many outgoing edges, the path space $X_B$ is locally compact. Lemma \ref{lemma_bdd_size_cone}  describes the structure of every generalized Bratteli diagram of bounded size. This result can be found in \cite[Section 3.1]{BezuglyiJorgensenKarpelSanadhya2025}.

\begin{lemma}\cite{BezuglyiJorgensenKarpelSanadhya2025}\label{lemma_bdd_size_cone} Let $B=(V,E)$ be a 
diagram of bounded size corresponding to a sequence 
$(t_n, L_n)_{n \in \N_0}$. Fix $v \in V_{n+1}$. Then 
for all infinite paths $x = 
(x_n)_{n\in \N_0}$ such that $r(x_n) = v$ and all $m > n$ we have
\begin{equation}\label{formula_bdd_size_lower_cone}
r(x_{m}) \in \left\{v - \sum_{i = n+1}^{m} t_{i}, \ldots, v + 
\sum_{i = n+1}^{m} t_{i}\right\},
\end{equation}
and for 
every $k \leq n$ we have 
$$
s(x_k) \subset \left\{v - \sum_{i = k}^{n} t_i, \ldots, v + 
\sum_{i = k}^n t_i\right\}.
$$
\end{lemma}

Let $B$ be a generalized Bratteli diagram of bounded size, $n \in \N_0$ and $v \in V_{n+1}$. We call the set of vertices $\{s(E(V_m, v)): m \leq n\}$ the \textit{upper cone} corresponding to $v$. The set of vertices $\{r(E(V_m, v)): m > n+1\}$ is called the the \textit{lower cone} corresponding to $v$.

For $w \in V_0$, define 
$$
Z_w^+ = \left\{x = (x_n) \in X_B : s(x_0) \geq w \mbox{ and } 
r(x_n) \geq w + \sum_{i = 0}^{n} t_i \mbox{ for all } n \in 
\mathbb{N}_0\right\}.
$$
and
$$
Z_w^- = \left\{x = (x_n) \in X_B : s(x_0) \leq w \mbox{ and } 
r(x_n) \leq w - \sum_{i = 0}^{n} t_i \mbox{ for all } n \in 
\mathbb{N}_0\right\}.
$$
The sets $Z^+_w$, $Z^-_w$ are called \textit{slanting sets} for 
$w \in V_0$. 
The following result can be found in \cite[Lemma 5.3, Theorem 5.4]{BezuglyiJorgensenKarpelSanadhya2025}.

%It was proved in \cite[Section 5]{BezuglyiJorgensenKarpelSanadhya2025} that for every $w \in V_0$, the slanting sets $Z_w^+$, $Z^-_w$ are closed nowhere dense sets which are invariant with respect to the tail equivalence relation. Hence, the tail equivalence relation for a generalized Bratteli diagram of bounded size is not minimal.

%\tcr{check extreme cases, e.g. $t_n^+ = 0$, don't as for the same $t_n^+$, $t_n^-$?}

\begin{theorem}\cite{BezuglyiJorgensenKarpelKwiatkowski2025}
    Let $B = (V,E) $ be a generalized Bratteli diagram of bounded size, with the 
corresponding sequence $(t_n, L_n)_{n \in \N_0}$. 
Then, for every $w \in V_0$, the sets $Z^+_w$, $Z^-_w$ are closed 
nowhere dense $\mathcal{R}$-invariant sets. In particular,  the tail equivalence relation 
$\mathcal{R}$ is \textit{not minimal}.
\end{theorem}

The tail equivalence relation for the Bratteli digrams of bounded size can be topologically transitive or not depending on the diagram. The following result was proved in \cite[Theorem 5.1]{BezuglyiJorgensenKarpelSanadhya2025}:

\begin{theorem}\label{ThmBJKStoptrans}\cite{BezuglyiJorgensenKarpelSanadhya2025}
 Let $B = (V,  E)$ be a generalized stationary Bratteli diagram 
with an 
irreducible aperiodic incidence matrix $F = (f_{ij})_{i,j \in 
\mathbb{Z}}$. Then the tail equivalence relation $\mathcal{R}$ is 
topologically transitive.   
\end{theorem}

It was also noticed in \cite[Remark 5.2]{BezuglyiJorgensenKarpelSanadhya2025} that if $B(V, E)$ is an irreducible generalized Bratteli diagram and
%with a vertex $i$ such that for every level $n$ there is an edge between $i \in V_n$ and $i \in V_{n+1}$ (in other words, 
there exists an infinite vertical path $x \in X_B$ 
%which passes through vertex $i$ on each level of the diagram) then
then the orbit $[x]_{\mathcal R}$ is dense in $X_B$, and 
the tail equivalence relation $\mathcal{R}$ is topologically 
transitive. Thus, we obtain the following corollary.

\begin{corol}
Let $B$ be a generalized Bratteli diagram of bounded size and $\mathcal{R}$ be the corresponding tail equivalence relation.

\begin{enumerate}[label=\upshape(\roman*), leftmargin=*, widest=iii]

\item If $B = B(F)$ is stationary and $F$ is irreducible then $\mathcal{R}$ is topologically transitive. In particular, if for some $t \in \N$ the incidence matrix $F$ has the property $f_{ij} > 0$ for $|i-j| \leq t$ then $F$ is irreducible and $\mathcal{R}$ is topologically transitive.

\item Let $B = B(F_n)$ have parameters $(t_n, L_n)$ such that
$f_{ij}^{(n)} > 0$ for all $|i - j| \leq t_n$. Then $B$ is irreducible, has vertical paths, and $\mathcal{R}$ is topologically transitive.
\end{enumerate}
\end{corol}

%Thus, for stationary generalized Bratteli diagrams of bounded size, the irreducibility of the incidence matrix implies topological transitivity of the tail equivalence relation. In particular, if for some $t \in \N$ the incidence matrix $F$ has the property $f_{ij} > 0$ for $|i-j| \leq t$ then $F$ is irreducible. In general, if the Bratteli diagram of bounded size has parameters $(t_n, L_n)$ and $f_{ij}^{(n)} > 0$ for all $|i - j| \leq t_n$ then the diagram is irreducible and has vertical paths, thus, the tail equivalence relation is topologically transitive.

\begin{example}[irreducible stationary diagram of bounded size with transitive tail equivalence relation]
The two-sided generalized Bratteli diagram $B(F)$ mentioned in Example \ref{Ex:isom_bd} and Remark \ref{Rem:isomvsirred} (see the upper diagram on Figure \ref{Fig:Isom_BD}) is a stationary irreducible generalized Bratteli diagram of bounded size with topologically transitive non-minimal tail equivalence relation. 
\end{example}


\begin{example}[irreducible stationary diagram of bounded size with non-transitive tail equivalence relation]\label{ex:bdd_irred_nontrans}
Let $B$ be a stationary generalized Bratteli diagram of bounded size with the incidence matrix $F = (f_{ij})$ defined as follows:
$$
f_{ij} = 
\left\{
\begin{aligned}
& 1, \mbox{ for } |i - j| = 1,\\
& 0, \mbox{ otherwise. }
\end{aligned}
\right.
$$
Let $\mathcal{R}$ be the tail equivalence relation on $B$.
Then $B$ is irreducible but $\mathcal{R}$ is not topologically transitive. Note that from even vertices on even levels of $B$ we can reach only odd vertices on odd levels of $B$, and vice versa, from odd vertices on even levels, we can reach only even vertices on odd levels. Thus, the diagram $B$ can be represented as a union of two disjoint subdiagrams $\ov B_{1}$ and $\ov B_2$, where the vertices of $\ov B_{1}$ are all even vertices on even levels and all odd vertices on odd levels, and the vertices of $\ov B_2$ are all odd vertices on even levels and all even vertices on odd levels. We have $X_B = X_{\ov B_{1}} \bigsqcup X_{\ov B_2}$, and the path spaces $X_{\ov B_{1}}$ and $X_{\ov B_2}$ are  $\mathcal{R}$-invariant. Since each $X_{\ov B_i}$, $i = 1,2$ is a countable union of cylinder sets of level $0$, both sets $X_{\ov B_{1}}$ and $X_{\ov B_2}$ are clopen.
Hence, the tail equivalence relation on $B$ is not topologically transitive. It is easy to see that the diagram is irreducible: from every vertex $i \in V_0$ (since the diagram $B$ is stationary, it's enough to consider only level $V_0$), we will eventually reach arbitrary vertex $j$ on some level below. The above phenomenon takes place because the matrix $F$ is not aperiodic: it has period $2$ (see e.g. \cite[Chapter 7]{Kitchens1998} or \cite[Appendix A]{BezuglyiJorgensenKarpelSanadhya2025}).
\end{example}

\begin{example}[reducible stationary diagram of bounded size with non-transitive tail equivalence relation]\label{ex:bdd_red_nontrans} We can modify Example \ref{ex:bdd_irred_nontrans} to obtain a reducible generalized Bratteli diagram $\tl B$ of bounded size with non-transitive tail equivalence relation. It suffices to choose $f_{ij} = 1$ for $|i-j| = 2$ and $f_{ij} = 0$ otherwise. Then from even vertices we can only reach even vertices, and from odd vertices we can only reach odd vertices. Thus, such diagram is reducible. The tail equivalence relation is non-transitive for the similar reason as in Example \ref{ex:bdd_irred_nontrans}: the diagram $\tl B$ can be represented as a disjoint union of four subdiagrams.    
\end{example}

%\tcr{do we need this example? No Does it split into disjoint subdiagrams? Yes}
% \begin{example}[reducible generalized Bratteli diagram of bounded size]
% Let $B  = (B,\omega)$ be a 
% generalized Bratteli diagram of bounded size, with the 
% corresponding sequence 
% $(t_n, L_n)_{n \in \N_0}$ defined as follows:  
% $$
% L_n = 2, \,\,\, t_n = 2^{n-1} \,\,\,\mathrm{for} \,\,\, n \in \N 
% \,\,\,\mathrm{and}\,\,\, t_0 = 1
% $$ 
% and vertices of each level identified with $\mathbb{Z}$. The diagram $B = (V,E)$ is shown in Figure \ref{Fig:ContinVMap}. 
% In other words, every vertex $v \in V \setminus V_0$ 
% has two incoming edges and every vertex $v \in V$ has two outgoing edges. It is easy to see that the diagram is reducible. For instance, since $t_n = 2^{n-1}$, if a vertex $v \in V_2$ is even then from $v$ we can only reach even vertices on lower levels. 

% %\tcr{What are the properties of the corresponding continuous Vershik map? Is it not transitive? No: diagram is split into two. Invariance of these properties under continuous/discountinous rernumerations of the vertices?}

%     \begin{figure}[ht]
% \unitlength=1cm
% \begin{graph}(11,8)
% % \graphnodesize{0.2}
% % \roundnode{V0}(3,6)
% %  %\nodetext{V0}(-1,0){$V_0$}
%  \roundnode{V11}(2,7)
%  %\nodetext{V11}(-0.7,0){$w_1^{(0)}$}
%   %\nodetext{V12}(0.7,0){$w_2^{(0)}$}
%  \roundnode{V12}(4,7)
%  \roundnode{V13}(6,7)
%  \roundnode{V14}(8,7)
%   \roundnode{V15}(10,7)


%   %
%     % The second level vertices
%    \roundnode{V21}(2,5)
%  %\nodetext{V11}(-0.7,0){$w_1^{(0)}$}
%   %\nodetext{V12}(0.7,0){$w_2^{(0)}$}
%  \roundnode{V22}(4,5)
%  \roundnode{V23}(6,5)
%  \roundnode{V24}(8,5)
%   \roundnode{V25}(10,5)
%   % The third level vertices
%  \roundnode{V31}(2,3)
%  \roundnode{V32}(4,3)
%  \roundnode{V33}(6,3)
%   \roundnode{V34}(8,3)
%     \roundnode{V35}(10,3)
  
%   % The fourth level vertices 
%   \roundnode{V41}(2,1)
%  \roundnode{V42}(4,1)
%  \roundnode{V43}(6,1)
%   \roundnode{V44}(8,1)
%     \roundnode{V45}(10,1)
  
%  %
%  % EDGES
%  \graphlinewidth{0.025}
% % % First level
% %  \edge{V0}{V11}
% %  \edge{V0}{V12}

%      \edge{V22}{V11}
%     \edge{V21}{V12}
     
%  \edge{V23}{V12}
%     \edge{V22}{V13}
    
%    \edge{V24}{V13}
%     \edge{V23}{V14}
    
%    \edge{V25}{V14}
%     \edge{V24}{V15}
    
%  % Second level
% \edge{V32}{V21}
%     \edge{V31}{V22}
     
%  \edge{V33}{V22}
%     \edge{V32}{V23}
    
%    \edge{V34}{V23}
%     \edge{V33}{V24}
    
%    \edge{V35}{V24}
%     \edge{V34}{V25}
    
%      % Third level
     
%  %\edge{V42}{V31}
%     \edge{V41}{V33}
     
%  \edge{V43}{V31}
%     \edge{V42}{V34}
    
%    \edge{V44}{V32}
%     \edge{V43}{V35}
    
%    \edge{V45}{V33}
%     %\edge{V44}{V35}
    
%   \freetext(10.9,7){$\ldots$}
%   \freetext(10.9,5){$\ldots$}
%  \freetext(10.9,3){$\ldots$}
%   \freetext(10.9,1){$\ldots$}
%    \freetext(1,7){$\ldots$}
%   \freetext(1,5){$\ldots$}  
%    \freetext(1,3){$\ldots$}
%   \freetext(1,1){$\ldots$}
%   \freetext(2,0.5){$\vdots$}
%   \freetext(4,0.5){$\vdots$}
%     \freetext(6,0.5){$\vdots$}
%       \freetext(8,0.5){$\vdots$}  
%         \freetext(10,0.5){$\vdots$} 
% \end{graph}
% \caption{A horizontally stationary reducible generalized Bratteli diagram}\label{Fig:ContinVMap}
% \end{figure} 
% \end{example}


\subsection{The renewal shift}\label{subsec:BD_renewal}

Consider the one-sided stationary generalized diagram $B_{RS} = B_{RS}(F_{RS})$, 
where the incidence matrix $F_{RS}$ is defined as follows:
\begin{equation}
    F_{RS} = \begin{pmatrix}
    1 & 1 & 0 & 0 & 0 & 0 & \cdots\\
    1 & 0 & 1 & 0 & 0 & 0 & \cdots\\
    1 & 0 & 0 & 1 & 0 & 0 & \cdots\\
    1 & 0 & 0 & 0 & 1 & 0 & \cdots\\
    1 & 0 & 0 & 0 & 0 & 1 & \cdots\\
    1 & 0 & 0 & 0 & 0 & 0 & \cdots\\
   \vdots & \vdots & \vdots & \vdots & \vdots & \vdots & \ddots
    \end{pmatrix}.
\end{equation}

\begin{figure}[hbt!]
\unitlength=1cm
\begin{graph}(9,4)
  %
    % The second level vertices
   \roundnode{V21}(0,3)
 %\nodetext{V11}(-0.7,0){$w_1^{(0)}$}
  %\nodetext{V12}(0.7,0){$w_2^{(0)}$}
 \roundnode{V22}(2,3)
 \roundnode{V23}(4,3)
 \roundnode{V24}(6,3)
  \roundnode{V25}(8,3)
  % The third level vertices
 \roundnode{V31}(0,1)
 \roundnode{V32}(2,1)
 \roundnode{V33}(4,1)
  \roundnode{V34}(6,1)
    \roundnode{V35}(8,1)
  
 %
 % EDGES
 \graphlinewidth{0.025}
% % First level
%  \edge{V0}{V11}
%  \edge{V0}{V12}



    
 % Second level

    \edge{V31}{V21}
        \edge{V31}{V22}
      
      \edge{V32}{V21} 
   \edge{V32}{V23}
     
 \edge{V33}{V21}
  \edge{V33}{V24}

    
   \edge{V34}{V21} 
    \edge{V34}{V25} 

   
   \edge{V35}{V21}

    
    
 \freetext(8.9,3){$\ldots$}
  \freetext(8.9,1){$\ldots$}
  
  \freetext(0,0.5){$\vdots$}
  \freetext(2,0.5){$\vdots$}
    \freetext(4,0.5){$\vdots$}
      \freetext(6,0.5){$\vdots$}  
        \freetext(8,0.5){$\vdots$} 


\end{graph}
\caption{One-sided stationary generalized Bratteli diagram $B_{RS}$ which corresponds to the renewal shift.}\label{Fig:BDrenewal}
\end{figure}


In other words, the vertex $1$ of each level is connected by a single edge to all vertices on the level below, and for $n > 1$, the vertex $n$ is connected by a singe edge to the vertex $(n-1)$ on the level below (see Figure \ref{Fig:BDrenewal}). 
The matrix $A_{RS} = F_{RS}^T$ appears in the theory of countable Markov shifts and corresponds to the renewal shift (see \cite[Chapter 7]{Kitchens1998} and \cite[Subsection 4.3.1]{Raszeja2021}, \cite[Subsection 3.2]{BEFR2022}). 
The diagram $B_{RS}$ was considered in \cite[Example 7.16]{BezuglyiJorgensenKarpelSanadhya2025}.
This diagram possesses a unique ergodic invariant probability measure which takes positive values on all cylinder sets (see \cite[Theorem 7.2, Proposition 7.17]{BezuglyiJorgensenKarpelSanadhya2025}). One can also find modifications of the diagram $B_{RS}$ and the corresponding incidence matrix, e.g., in \cite[Subection 7.3]{BezuglyiJorgensenKarpelSanadhya2025} and \cite[Subection 4.3]{Raszeja2021}.
%It was proved in \cite[Proposition 7.17]{BezuglyiJorgensenKarpelSanadhya2025} that the generalized Bratteli diagram corresponding to the renewal shift supports a unique probability ergodic tail invariant measure. 
%$\mu$ defined using the Perron eigenvalue $\lambda = 2$ and the corresponding right eigenvector $\xi = \left(\frac{1}{2}, \frac{1}{2^2}, \frac{1}{2^3}, \ldots\right)^T$ of $A_{BS}$. 
%The matrix $A_5$ is positive recurrent.




Theorem \ref{thm:minimalRrenewalshift} below shows that the diagram $B_{RS}$ provides an example of an irreducible generalized Bratteli diagram with a minimal tail equivalence relation. 
%This result formulated in terms of a Vershik map was mentioned without a proof in \cite[Example 4.2]{BezuglyiJorgensenKarpelSanadhya2025} for a similar Bratteli diagram. 
%Here, for reader's convenience, we give a detailed proof for the diagram $B_{RS}$.

\begin{theorem}\label{thm:minimalRrenewalshift}
    Let $\mathcal{R}_{RS}$ be the tail equivalence relation on the generalized Bratteli diagram $B_{RS}$ which corresponds to the renewal shift. Then $\mathcal{R}_{RS}$ is minimal and $B_{RS}$ is irreducible.
\end{theorem} 

\begin{proof} First, we show that $B_{RS}$ is irreducible.
From the vertex $i = 1$ on any level of the diagram $B_{RS}$, there are edges to any vertex $j \in \N$ on a level below. 
For every $n \in \N_0$ and any vertex $i > 1$ in $V_n$, there is a finite path from $i$ to the first vertex on the level $n + i - 1$. Hence, from the vertex $i$ on the level $n$, we can reach any vertex $j$ on the level $n + i$. Thus, the diagram $B_{RS}$ is irreducible.

Now we show that $\mathcal{R}_{RS}$ is minimal. Due to the structure of $B_{RS}$ (see Figure \ref{Fig:BDrenewal}), any path from $X_{B_{RS}}$ passes through the first vertex infinitely many times. Let $x \in X_{B_{RS}}$ be an arbitrary path and $[(y_0, \ldots, y_n)]$ be an arbitrary cylinder set. Then $r(y_n) \in V_{n+1}$ is joined by a finite path to the first vertex on some level $m > n+1$. 
%(if $r(y_n) = 1$ set $m = n+1$). 
Since $x$ passes through the first vertex infinitely many times, there is a level $m_1 > m$ such that $s(x_{m_1}) = 1$. There is a vertical path in $B$ between the first vertex on level $m$ and the first vertex on level $m_1$, hence, there is an infinite path $z = (z_i)$ such that $z_i = x_i$ for $i \geq m_1$, $s(z_i) = 1$ for $m \leq i \leq m_1$ and $z_i = y_i$ for $0 \leq i \leq n$. Thus, $z$ is tail equivalent to $x$ and belongs to the set $[(y_0, \ldots, y_n)]$. In other words, the orbit of $x$ under the tail equivalence relation visits $[(y_0, \ldots, y_n)]$. Thus, the orbit
of every infinite path under the tail equivalence relation visits every cylinder set, thus the tail equivalence relation $\mathcal{R}_{RS}$ is minimal.  
\end{proof}



\subsection{Diagrams of infinite odometers}\label{subsec:DIO}
A class of generalized Bratteli diagrams called \textit{diagrams of infinite odometers} ($B_{IO}$) was introduced in \cite[Section 4]{BezuglyiKarpelKwiatkowski2024}. It's incidence matrices are $\N \times \N$ upper triangular matrices
$$
{F}_{n} =
\begin{pmatrix}
 a_{1}^{(n)} &     1 &      0 &      0 & \ldots  \\
     0 &  a_{2}^{(n)} &      1 &      0 &      \ldots  \\
     0 &      0 &  a_{3}^{(n)} &      1 &      \ldots  \\
     0 &      0 &      0 &  a_{4}^{(n)} &      \ldots  \\
\vdots & \vdots & \vdots & \vdots & \ddots \\
\end{pmatrix}, \quad n \in \N_0,
$$
with entries
$$
f_{ij}^{(n)} = 
\left\{
\begin{aligned}
& a_{i}^{(n)} \geq 2, \mbox{ for } j = i,\\
& 1, \mbox{ for } j = i + 1,\\
& 0, \mbox{ otherwise. }
\end{aligned}
\right.
$$


By \textit{odometer} we will mean a standard Bratteli diagram which has a single vertex on each level and at least two edges connecting the vertices of neighbouring levels (see \cite[Subsection 6.5.1]{Durand2010}). For $i \in \N$, denote by $B_i$ the vertical odometer which passes through the vertex $i$ on each level. In other words, odometer $B_i$ has $a_n^{(i)}$ edges connecting the vertex of level $n$ with the vertex of level $n + 1$. Denote by $X_{B_i}$ the set of all paths from $X_B$ that pass through vertex $i$ on each level. Denote by $\wh {X}_{B_i}$ the set of all paths from $X_B$ that are tail equivalent to a path from $X_{B_i}$. Then $X_{B_{IO}} = \bigsqcup_{i = 1}^{\infty}\wh X_{\ov B_i}$ (see \cite[Proposition 4.1]{BezuglyiKarpelKwiatkowski2024}). In other words, every infinite path is eventually vertical and is tail equivalent to a path from $B_i$ for some $i \in \N$.


For diagram $B_{IO}$, one can characterize the set of all ergodic invariant measures (see \cite[Theorem 4.3]{BezuglyiKarpelKwiatkowski2024}). In particular, diagrams of infinite odometers provide examples of generalized Bratteli diagrams which possess no finite ergodic invariant measures. For instance, this happens for $a_{i}^{(n)} = i + 1$ for all $n \in \N_0$, see Figure \ref{Fig:no measure}.

\begin{figure}[hbt!]
\unitlength=.8cm
\begin{graph}(7,6)
% \graphnodesize{0.2}
% \roundnode{V0}(3,6)
%  %\nodetext{V0}(-1,0){$V_0$}
 \roundnode{V11}(2,5)
 %\nodetext{V11}(-0.7,0){$w_1^{(0)}$}
  %\nodetext{V12}(0.7,0){$w_2^{(0)}$}
 \roundnode{V12}(4,5)
 \roundnode{V13}(6,5) 
  % The second level vertices
 \roundnode{V21}(2,3)
 \roundnode{V22}(4,3)
 \roundnode{V23}(6,3)
  %\nodetext{V21}(-0.7,0){$w_1^{(1)}$}
 % \nodetext{V22}(0.7,0){$w_2^{(1)}$}
 % The third level vertices
 \roundnode{V31}(2,1)
 \roundnode{V32}(4,1)
 \roundnode{V33}(6,1)
 % \nodetext{V31}(-0.7,0){$w_1^{(2)}$}
  %\nodetext{V32}(0.7,0){$w_2^{(2)}$}
  %
 %
 % EDGES
 \graphlinewidth{0.025}
% % First level
%  \edge{V0}{V11}
%  \edge{V0}{V12}

 % Second level
 \bow{V21}{V11}{0.09}
  \bow{V21}{V11}{-0.09}
    \edge{V21}{V12}
     
 \bow{V22}{V12}{0.09}
 \bow{V22}{V12}{-0.09}
 \edge{V22}{V12}
 
 \bow{V23}{V13}{0.12}
  \bow{V23}{V13}{-0.12}
   \bow{V23}{V13}{0.04}
  \bow{V23}{V13}{-0.04}
 \edge{V22}{V13}

 %third level
 \bow{V31}{V21}{0.09}
 \bow{V31}{V21}{-0.09}
   \edge{V31}{V22}
%     \edge{V32}{V22}[\graphlinecolour(1,0,0)]
 \bow{V32}{V22}{0.09}
 \bow{V32}{V22}{-0.09}
 \edge{V32}{V22}
 
 
  \bow{V33}{V23}{0.12}
  \bow{V33}{V23}{-0.12}
   \bow{V33}{V23}{0.04}
  \bow{V33}{V23}{-0.04}
 \edge{V32}{V23}
 
    \freetext(6.9,5){$\ldots$}
  \freetext(6.9,3){$\ldots$}
    \freetext(6.9,1){$\ldots$}
    \freetext(2,0.1){$\vdots$}
  \freetext(4,0.1){$\vdots$}
    \freetext(6,0.1){$\vdots$}
%\freetext(3,0.1){.\,.\,.\,.\,.\,.\,.\,.\,.\,.\,.\,.\,.\,.\,.\,.\,.\,.\,.\,.\,..\,.\,.\,.\,.\,.\,.\,.}
\end{graph}
\caption{A stationary one-sided $B_{IO}$ that does not have finite invariant measures.}\label{Fig:no measure}
\end{figure}

In this paper, we will call diagrams $B_{IO}$ \textit{one-sided diagrams of infinite odometers}. We also introduce \textit{two-sided Bratteli diagrams of infinite odometers} $\tl B_{IO}$ with $\Z \times \Z$ incidence matrices $\tl F_n$ that have entries $a_i^{(n)} \geq 2$ on the main diagonal and ones on the diagonal above:
    $$
    \tl F_n = 
    \begin{pmatrix}
       \ddots & \vdots & \vdots & \vdots & \vdots & \udots\\
       \ldots & \textbf{$a_{-1}^{(n)}$} & 1 & 0 & 0 &\ldots\\
       \ldots & {0} & \textbf{$a_0^{(n)}$} & 1 & 0 & \ldots\\   
       \ldots & {0} & 0 & \textbf{$a_1^{(n)}$} & 1 & \ldots\\ 
       \ldots & 0 & 0 & 0 & \textbf{$a_{2}^{(n)}$} & \ldots\\     
       \udots & \vdots & \vdots & \vdots & \vdots & \ddots
    \end{pmatrix}.
    $$

%Clearly, a one-sided diagram $B_{IO}$ is reducible since from every vertex $i$ we cannot reach any vertex $j > i$ on any level below. It is also easy to see that the tail equivalence relation on $B_{IO}$ is topologically transitive. Indeed, since from any vertex we can reach the first vertex on some level below, any path that passes through the first vertex on each level of the diagram has a dense orbit under the tail equivalence relation.

The structure of the two-sided diagram $\tl B_{IO}$ is more complicated than the one of $B_{IO}$ since not every path in $X_{\tl B_{IO}}$ is eventually vertical. Since $\tl F_n$ has only values $1$ right above the main diagonal, the set of slanted paths is countable (see also \cite[Proposition 5.5]{BezuglyiJorgensenKarpelSanadhya2025}).

\begin{prop}
Both one-sided and two-sided Bratteli diagrams of infinite odometers $B_{IO}$ and $\tl B_{IO}$ are reducible, their corresponding tail equivalence relations $\mathcal{R}_{IO}$ and $\widetilde{\mathcal{R}}_{IO}$ are topologically transitive but not minimal. 
\end{prop}

\begin{proof}
For both diagrams $B_{IO}$ and $\tl B_{IO}$, from a vertex $i$ on level $n$ we will never reach vertex $j > i$ on any level below. Hence, both diagrams $B_{IO}$ and $\tl B_{IO}$ are reducible.

For the one-sided diagram $B_{IO}$, 
the set $\wh{X}_{B_{RS}}$ is dense in $X_{B_{IO}}$. The orbit a path $x \in X_{B_{IO}}$ is dense in $X_{B_{IO}}$ if and only if it $x$ is tail equivalent to a path from $X_{B_{RS}}$. Thus, the tail equivalence relation $\mathcal{R}_{IO}$ is topologically transitive but not minimal.

%any vertical path which passes through the vertex $0$ has a dense orbit, but the orbit of a vertical path which passes through the vertex $1$ will never visit the cylinder $[0]$ which corresponds to the vertex $0 \in V_0$.

For the two-sided diagram of infinite odometers $\tl B_{IO}$, there is no $i \in \Z$ such that $\wh{X}_{B_i}$ is dense in $X_{\tl B_{IO}}$. The tail equivalence class of any eventually vertical path is not dense, but we can find other paths that have dense orbits under the tail equivalence relation on the path space of $\tl B_{IO}$. 
First, note that any path $x \in X_{\tl B_{IO}}$ divides the set of all vertices of $\tl B_{IO}$ into three disjoint nonempty subsets: the set $V(x) = \{s(x_n)\}_{n = 0}^{\infty}$ of vertices through which $x$ passes, the set 
$$V_{l}(x) = \bigcup_{n \geq 0}\{w \in V_n : w < s(x_n)\}$$ 
%of vertices that are above or to the left of $x$, 
%such that $w \in V_n$ belongs to $V_l(x)$ if $w < s(x_n)$,
and the set 
$$V_{r}(x) = \bigcup_{n \geq 0}\{w \in V_n : w > s(x_n)\}.$$ 
%$V_r(x)$ 
%of the vertices %which are below or to the right of $x$.

Clearly, a path with a dense orbit cannot be vertical. Indeed, the paths from the tail equivalence class of a vertical path passing through the vertex $i$ on each level can reach only vertices $j \geq i$ on the levels above, i.e. the vertices from $V_r(x)$. Thus, they will never belong to a cylinder set which ends in a vertex from $V_l(x)$. In particular, they will never visit a cylinder set which corresponds to the vertex $i - 1$ on level $0$. 
     The slanted paths which start at some vertex $i \in V_0 = \Z$ and go all the time through the slanted edges also don't have dense orbits. Indeed, orbit of such a path will never reach a cylinder set which corresponds to the vertex $i+1$ on level $0$ or, in general, any other vertex from $V_r(x)$.
     
     The path $x$ which starts at a vertex $i \in V_0$ and passes in turn through vertical and slanted edges has a dense orbit. In other words, we consider the path which passes through the vertices $i, i, i - 1, i - 1, \ldots, i - k, i - k, \ldots$. 
     For simplicity, set $i = 0$, then we have $s(x_{2n}) = s(x_{2n + 1}) = -n$. Thus, $x$ passes through the vertex number $-n$ on levels $2n+1$ and $2n+2$. 
     %Note that the tower which corresponds to the vertex $i$ on level $l$ reaches the vertices $i, \ldots, i + (l - k)$ on the level $k < l$. Thus, the tower which corresponds to vertex $-n$ on the level $l = 2n$, reaches the vertices $-n, \ldots \tcr{compute}$
      Let $[(y_0, \ldots, y_n)]$ be any cylinder set that ends at the vertex $j^{(n+1)}$ on the level $n + 1$. If $j^{(n+1)} \in V_{l}(x)$, then we can use a vertical path that starts at $j^{(n+1)}$ to reach a vertex from $V(x)$. Thus, a point from the orbit of $x$ will belong to the cylinder $[(y_0, \ldots, y_m)]$. If $j^{(m+1)} \in V_{r}(x)$, we use the slanted path $sl(j^{(m+1)})$ that starts at $j^{(m+1)}$ to reach a vertex from $V(x)$. Such a path passes through the vertex $j^{(m+1)} - k$ on the level $m + 1 + k$.  If $j^{(m+1)} > 0$ then the path $sl(j^{(m+1)})$ passes through the vertex $0$ on the level $m + 1 + j^{(m+1)}$ and it passes though vertex $-n$ on the level $m + 1 + j^{(m+1)} + n$. Thus, for $n = j^{(m+1)} + m$ and $n = j^{(m+1)} + m - 1$, the slanted path $sl(j^{(m+1)})$ passes through the vertex which belongs to $V(x)$. If $j^{(m+1)} \leq 0$, we shift all the vertices uniformly by the same number $-r(x_m)$ to the right, so that now we have $r(x_m) = 0$, and consider only the part of $\tl B_{IO}$ which starts at the level $m+1$, Thus, we repeat the same proof as in the previous case, i.e. the proof for $j^{(m+1)} > 0$.  

      Note that the same proof works for any path $x \in X_{\tl B_{IO}}$ has infinitely many vertical and infinitely many slanted edges. In other words, $x \in X_{\tl B_{IO}}$ has a dense orbit if and only if $x$ is not tail equivalent to a vertical or to a slanted path. Such a point $x$ will ``intersect'' any slanted path that starts at a vertex from $V_r(x)$ and any vertical path that starts at a vertex from $V_{l}(x)$. Since $\tl B_{IO}$ contains only edges which are vertical or slanted that join a vertex $i$ to a vertex $i - 1$ on a level below, any path which passes through vertices from both $V_r(x)$ and $V_l(x)$ should pass through a vertex from $V(x)$. 
      %one cannot move from V_r on level n to V_l on level n+1.  Or from V_l above to V_r below.
      %some geometry here
      %for more explanation consider cases: V_l(x), V(x), V_r(x) on a level above and below, and if x has a vertical or a slanted edge 
\end{proof}


\subsection{Diagram $B_{\infty}$}\label{subsec:diagramB-infty}
Diagram $B_{\infty}$ is a one-sided generalized Bratteli diagram defined by the lower triangular $\N \times \N$ incidence matrix $F_{\infty}$ which has all entries equal to $1$ on the main diagonal and below the main diagonal (see Figure \ref{Fig:Binfty}):
% $$
% f_{ij} = 
% \left\{
% \begin{aligned}
% & 1, \mbox{ for } i \geq j,\\
% & 0, \mbox{ otherwise. }
% \end{aligned}
% \right.
% $$

$$
F_{\infty} = 
\begin{pmatrix}
    1 & 0 & 0 & 0 & \ldots\\
    1 & 1 & 0 & 0 & \ldots\\
    1 & 1 & 1 & 0 & \ldots\\
    1 & 1 & 1 & 1 & \ldots\\
    \vdots & \vdots & \vdots & \vdots & \ddots
    \end{pmatrix}
$$



\begin{figure}[hbt!]
\unitlength=1cm
\begin{graph}(9,4)
  %
    % The second level vertices
   \roundnode{V21}(0,3)
 %\nodetext{V11}(-0.7,0){$w_1^{(0)}$}
  %\nodetext{V12}(0.7,0){$w_2^{(0)}$}
 \roundnode{V22}(2,3)
 \roundnode{V23}(4,3)
 \roundnode{V24}(6,3)
  \roundnode{V25}(8,3)
  % The third level vertices
 \roundnode{V31}(0,1)
 \roundnode{V32}(2,1)
 \roundnode{V33}(4,1)
  \roundnode{V34}(6,1)
    \roundnode{V35}(8,1)
  
 %
 % EDGES
 \graphlinewidth{0.025}
% % First level
%  \edge{V0}{V11}
%  \edge{V0}{V12}



    
 % Second level

    \edge{V31}{V21}
    
    \edge{V32}{V21}
        \edge{V32}{V22}
     
 \edge{V33}{V21}
  \edge{V33}{V22}
 \edge{V33}{V23}
    
   \edge{V34}{V21} 
    \edge{V34}{V22} 
       \edge{V34}{V23} 
   \edge{V34}{V24} 
   
   \edge{V35}{V21}
      \edge{V35}{V22}
         \edge{V35}{V23}
            \edge{V35}{V24}
   \edge{V35}{V25}
    
    
 \freetext(8.9,3){$\ldots$}
  \freetext(8.9,1){$\ldots$}
  
  \freetext(0,0.5){$\vdots$}
  \freetext(2,0.5){$\vdots$}
    \freetext(4,0.5){$\vdots$}
      \freetext(6,0.5){$\vdots$}  
        \freetext(8,0.5){$\vdots$} 


\end{graph}
\caption{One-sided stationary generalized Bratteil diagram $B_{\infty}$.}\label{Fig:Binfty}
\end{figure}

\smallskip

The diagram $B_{\infty}$ was introduced in \cite[Section 8]{BezuglyiKarpelKwiatkowskiWata2024} as an example of a diagram which has uncountably many probability ergodic invariant measures, and one can explicitly describe these measures using the eigenvectors and eigenvalues of the incidence matrix (see \cite[Theorem 4.3]{BezuglyiKarpelKwiatkowski2024}). The diagram is reducible, since from every vertex number $i > 1$ on a level $n \in \N_0$, we cannot reach a vertex number $j < i$ on a lower level. Proposition \ref{prop:B_infty} shows that the tail equivalence relation on $X_{B_{\infty}}$ is topologically transitive but not minimal. 
%It is easy to see that the orbit of any vertical path is not dense, since from a vertex $i$ on a level $n \in \N$ we can reach only vertices with numbers $j \leq i$ on any level above. More generally,
\begin{prop}\label{prop:B_infty}
Let $x = (x_n)$ be a path in $X_{B_{\infty}}$. Then the orbit $\mathcal{O}(x)$ is dense in $X_B$ if and only if the set $\{s(x_n)\}_{i = 0}^{\infty} \subset \N$ is unbounded.
\end{prop}

\begin{proof}
Fist, assume that we have $s(x_n) < M$ for some $M \in \N$ and all $n \in \N$. Notice that from a vertex $i$ on a level $n \in \N$ we can reach only vertices with numbers $j \leq i$ on any level above. Hence, $\mathcal{O}(x) \cap [(y_0, \ldots, y_n)] = \emptyset$ for every cylinder set $[(y_0, \ldots, y_n)]$ with $r(y_n) \geq M$.

Now assume that for every $M \in \N$ there exists $n \in \N$ such that $s(x_n) > M$. 
It follows from the structure of the diagram that $s(x_m) \geq s(x_n) > M$ for all $m \geq n$. Let $[(y_0, \ldots, y_n)]$ be a cylinder set. Then there exists $m \in \N$ such that $s(x_m) > r(y_n)$ and $m > n + 1$. Thus, from a vertex $s(x_m) \in V_m$, we can reach a vertex number $s(x_m)$ on level $V_{n+2}$, and since $s(x_m) > r(y_n)$, there is an edge between the vertex $r(y_n) \in V_{n+1}$ and the vertex number $s(x_m)$ on level $V_{n+2}$. Thus, there is a finite path between $r(y_n)$ and $s(x_m)$, and $\mathcal{O}(x) \cap [(y_0, \ldots, y_n)] \neq \emptyset$.
\end{proof}

%An example of a path with a dense orbit is the path which starts at vertex $1 \in V_0$ and passes through slanted edges on even levels $E_{2n}$ and vertical edges on odd levels $E_{2n+1}$ has a dense orbit under the tail equivalence relation. \tcr{complete characterisation of paths with dense orbits}
%\tcr{top transitive? What about slanted paths? Or exchanging slanted and vertical infinitely many times?}

% \begin{example}\label{Ex:reducible_GBD_1}

% Let $B$ be a stationary generalized Bratteli diagram with the set of vertices on each level identified with $\mathbb{N}$ and the entries of the incidence matrix $F$ given by the formula 
% $$
% f_{ij} = 
% \left\{
% \begin{aligned}
% & 1, \mbox{ for } i \geq j,\\
% & 0, \mbox{ otherwise. }
% \end{aligned}
% \right.
% $$
% Then we have
%     $$F = 
%     \begin{pmatrix}
%      1 & 0 & 0 & 0 &\ldots\\
%      1 & 1 & 0 & 0 &\ldots\\
%      1 & 1 & 1  & 0 &\ldots\\
%      1 & 1 & 1  & 1 & \ldots\\
%      \vdots & \vdots & \vdots & \vdots & \ddots
%     \end{pmatrix},
%     $$
% and for every $n = 0,1,2, \ldots$, the first vertex on level $n$ has an edge to every vertex on the level $n+1$. It is easy to see that diagram $B$ is reducible. Indeed, for every vertex $w > 1$ on any level $n$ there is no path to the first vertex on levels $m > n$.

% \begin{figure}[hbt!]
% \unitlength=1cm
% \begin{graph}(9,4)
%   %
%     % The second level vertices
%    \roundnode{V21}(0,3)
%  %\nodetext{V11}(-0.7,0){$w_1^{(0)}$}
%   %\nodetext{V12}(0.7,0){$w_2^{(0)}$}
%  \roundnode{V22}(2,3)
%  \roundnode{V23}(4,3)
%  \roundnode{V24}(6,3)
%   \roundnode{V25}(8,3)
%   % The third level vertices
%  \roundnode{V31}(0,1)
%  \roundnode{V32}(2,1)
%  \roundnode{V33}(4,1)
%   \roundnode{V34}(6,1)
%     \roundnode{V35}(8,1)
  
%  %
%  % EDGES
%  \graphlinewidth{0.025}
% % % First level
% %  \edge{V0}{V11}
% %  \edge{V0}{V12}



    
%  % Second level

%     \edge{V31}{V21}
    
%     \edge{V32}{V21}
%         \edge{V32}{V22}
     
%  \edge{V33}{V21}
%   \edge{V33}{V22}
%  \edge{V33}{V23}
    
%    \edge{V34}{V21} 
%     \edge{V34}{V22} 
%        \edge{V34}{V23} 
%    \edge{V34}{V24} 
   
%    \edge{V35}{V21}
%       \edge{V35}{V22}
%          \edge{V35}{V23}
%             \edge{V35}{V24}
%    \edge{V35}{V25}
    
    
%  \freetext(8.9,3){$\ldots$}
%   \freetext(8.9,1){$\ldots$}
  
%   \freetext(0,0.5){$\vdots$}
%   \freetext(2,0.5){$\vdots$}
%     \freetext(4,0.5){$\vdots$}
%       \freetext(6,0.5){$\vdots$}  
%         \freetext(8,0.5){$\vdots$} 


% \end{graph}
% \caption{}\label{Fig:ReducibleGBD}
% \end{figure}

% \end{example}


\section{Isomorphic generalized Bratteli diagrams}\label{sec:mainres}


\subsection{Properties invariant under isomorphism}

Let $B$ and $B'$ be (standard or generalized) Bratteli diagrams with the corresponding tail equivalence relations $\mathcal{R}$ and $\mathcal{R}'$.
Assume that $B$ is isomorphic to $B'$ via the sequences of vertex bijections $(g_n)_{n = 0}^{\infty}$ and edge bijections $(h_n)_{n = 0}^{\infty}$. It is straightforward to check that the map $\theta: X_B \rightarrow X_{B'}$ defined by $$\theta((x_n)_{n = 0}^{\infty}) = (h_n(x_n))_{n = 0}^{\infty}$$ is a homeomorphism which preserves tail equivalence relation. In other words, for all $x,y \in X_B$, we have $(x,y) \in \mathcal{R}$ if and only if $(\theta(x), \theta(y)) \in \mathcal{R}'$. Thus, all topological properties of the tail equivalence relation such as minimality or topological transitivity 
are preserved under the isomorphism of (standard or generalized) Bratteli diagrams. It also follows that $\mathcal{R}$ and $\mathcal{R}'$ are Borel isomorphic (see \cite{Kechris2024} for more details).

A matrix $F = (f_{ij})$ satisfies the \textit{equal row sum} property (we write $F\in ERS$ or $F\in ERS(r)$) if there exists $r$ such that $\sum_{j} f_{ij} = r$ for all $i$. A matrix $F = (f_{ij})$ has the {\em equal column sum} property ($F \in ECS$ or $F\in ECS(c)$) if there exists $c$ such that $\sum_{i} f_{ij} = c$ for all $j$. We say that a Bratteli diagram $B$ has the $ERS(r_n)$ ($ECS(c_n)$) property, if for all $n$ the incidence matrix $F_n \in ERS(r_n)$ ($ECS(c_n)$). Since isomorphism of Bratteli diagrams preserves levels, it preserves the properties of incidence matrices to have equal row (equal column) sum.

Hence, we obtain the following Proposition \ref{Prop:prop_pres_under_isom}.

\begin{prop}\label{Prop:prop_pres_under_isom}
    Let $B$ and $B'$ be two isomorphic (standard or generalized) Bratteli diagrams and $\mathcal{R}$, $\mathcal{R}'$ be the corresponding tail equivalence relations. Then
\begin{enumerate}[label=\upshape(\roman*), leftmargin=*, widest=iii]

\item If $B$ has the $ERS(r_n)$, $n \geq 0$ ($ECS(c_n)$, $n \geq 0$) property then $\tl B$ has the $ERS(r_n)$, $n \geq 0$ ($ECS(c_{n})$, $n \geq 0$) property.

\item  If $\mathcal{R}$ is topologically transitive (minimal) then $\mathcal{R'}$ is topologically transitive (minimal).

\end{enumerate}
\end{prop}

\begin{definition}
    A generalized Bratteli diagram $B$ is called \textit{connected} if it is connected as an undirected graph.
\end{definition}

Clearly, connectedness as an undirected graph is also preserved under isomorphism. The properties of being stationary or of bounded size are not preserved under isomorphism. Below we show that, in general, irreducibility is also not preserved under isomorphism. 

\subsection{Irreducibility and isomorphism}

%\tcr{say smth about standard diagrams}

%Till the end of this subsection, we will consider only generalized Bratteli diagrams. 
In this subsection, we  show that the set of all generalized Bratteli diagrams can be partitioned into two disjoint subsets:

\begin{enumerate} [label=\upshape(\roman*), leftmargin=*, widest=iii]
    \item The diagrams which can be isomorphic only to irreducible generalized Bratteli diagrams. We  call such diagrams \textit{completely irreducible}.

    
    \item The diagrams which can be isomorphic both to reducible and to irreducible  generalized Bratteli diagrams. We call such diagrams \textit{relatively irreducible}.
    
    %\item The diagrams which can be Borel isomorphic only to reducible generalized Bratteli diagrams (we will call them diagrams of Class III); - no such BD exist

\end{enumerate}

Theorem \ref{thm:classII} below shows that there are no generalized Bratteli diagrams that are isomorphic only to reducible generalized Bratteli diagrams. It other words, the vertices of every generalized Bratteli diagram can be enumerated in such a way that the diagram becomes an irreducible generalized Bratteli diagram.

First, note that the following result  holds.

\begin{prop}
\label{prop:toptransirred}
     Let $B$ be a generalized Bratteli diagram and $\mathcal{R}$ be the tail equivalence relation on $B$. If $\mathcal{R}$ is topologically transitive then $B$ is isomorphic to an irreducible generalized Bratteli diagram.
     %, hence $B$ cannot belong to Class III.
\end{prop}

\begin{proof}
    Let $x = (x_n) \in X_B$ be a point with dense orbit. Consider a diagram $B'=(V',E')$ isomorphic to $B$ such that the image $\psi(x)$ of $x$ under the isomorphism passes through the vertex number $n$ infinitely many times for every $n$. For instance, the vertices $g(s(x_n))$ might form the following sequence: $0, 0, 1, 0, 1, 2, 0, 1, 2, 3,\ldots$ Since $\mathcal{O}(x)$ is dense in $X_B$, for every cylinder set $[(e_0, \ldots, e_n)] \subset X_B$ there is a finite path between $r(e_n)$ and a vertex $s(x_m)$ for some $m$. Hence, for every cylinder set $[(e_0', \ldots, e_n')] \subset X_{B'}$ there is a finite path from $r'(e_n')$ to $s'(h(x_m)) = g_m(s(x_m))$ for some $m$. Since any number $j \in \N_0$ appears among $\{g_k(s(x_k))\}_{k \in \N_0}$ infinitely many times, we will reach any vertex $j$ from $g_m(s(x_m))$. Thus, the diagram $B'$ is irreducible.
    %\tcb{Conversely, assume that $\mathcal{R}$ is not topologically transitive, but $B$ is isomorphic to some irreducible generalized Bratteli diagram $B' = (V', E')$. Fix any vertex $j \in V_0'$. Since $B'$ is irreducible, there is a path from $j \in V_0'$ to $j \in V'_{n_1}$ for some $n_1 > 0$. Similarly, there is a path from $j \in V'_{n_1}$ to $j \in V'_{n_2}$ for some $n_2 > n_1$ and so on. Set $n_0 = 0$. Thus, we obtain an increasing sequence numbers $\{n_k\}_{k = 0}^{\infty} \subset \N_0$ such that there is a finite path in $B'$ between $j \in V_{n_k}$ and $j \in V_{n_{k+1}}$ for all $k$. Telescope \tcr{add def in Sect 2} $B'$ with respect to the levels $n_k$, and obtain a new diagram $B''$. Now, $B''$ has a vertical path passing through the vertex $j$ on each level.} - doesn't really work
\end{proof}

%The converse to Lemma \ref{thm:classIII} is not true (see Example below).

For instance, diagram $B_{\infty}$ (see Subsection \ref{subsec:diagramB-infty}) and diagrams of infinite odometers (see Subsection \ref{subsec:DIO}) have topologically transitive tail equivalence relations. Below we give an example of a generalized Bratteli diagram with a non topologically transitive tail equivalence relation. Moreover, this diagram 
has countably infinitely many distinct sets $\ov{\mathcal{O}(x^i)}$ such that $X_B = \bigcup_{i = 1}^{\infty} \ov{\mathcal{O}(x^i)}$.

\begin{example}\label{ex:connected_non-transR}
The one-sided stationary generalized Bratteli diagram $B$ shown on Figure \ref{Fig:nontransitnonloccomp} with vertices of each level enumerated by $\N$, has the incidence matrix $F = (f_{ij})$ defined as follows: $f_{11} = 2$; $f_{ii} = 3$ for $i > 1$; $f_{i1} = 1$ for $i > 1$; and $f_{ij} = 0$ otherwise.
% $$
% f_{ij} = 
% \left\{
% \begin{aligned}
% & 2, \mbox{ for } i = j = 1,\\
% & 3, \mbox{ for } i = j > 1,\\
% & 1, \mbox{ for } 1 = i < j,\\
% &0 , \mbox{ otherwise.}
% \end{aligned}
% \right.
% $$
%the first vertex of each level has the first consists of vertical odometers and 
The diagram has a unique minimal component for the tail equivalence relation: it is the dyadic odometer which passes through the first vertex on each level of the diagram. For $i \in \N$, let $x^i$ be a path in $X_B$ which eventually passes through the vertex number $i$ on each level of the diagram. Then $X_B = \bigcup_{i = 1}^{\infty} \ov{\mathcal{O}(x^i)}$ and $\ov{\mathcal{O}(x^i)} \cap \ov{\mathcal{O}(x^j)} = \ov{\mathcal{O}(x^1)}$ for all $i \neq j$. Note that all paths in $X_B$ are eventually vertical, and there is no path in $X_B$ with a dense orbit. Thus, the tail equivalence relation on $X_B$ is not topologically transitive.


\begin{figure}
    \unitlength=.8cm
\begin{graph}(9,6)
% \graphnodesize{0.2}
% \roundnode{V0}(3,6)
%  %\nodetext{V0}(-1,0){$V_0$}
 \roundnode{V11}(2,5)
 %\nodetext{V11}(-0.7,0){$w_1^{(0)}$}
  %\nodetext{V12}(0.7,0){$w_2^{(0)}$}
 \roundnode{V12}(4,5)
 \roundnode{V13}(6,5) 
  \roundnode{V14}(8,5)
  % The second level vertices
 \roundnode{V21}(2,3)
 \roundnode{V22}(4,3)
 \roundnode{V23}(6,3)
  \roundnode{V24}(8,3)
  %\nodetext{V21}(-0.7,0){$w_1^{(1)}$}
 % \nodetext{V22}(0.7,0){$w_2^{(1)}$}
 % The third level vertices
 \roundnode{V31}(2,1)
 \roundnode{V32}(4,1)
 \roundnode{V33}(6,1)
  \roundnode{V34}(8,1)
 % \nodetext{V31}(-0.7,0){$w_1^{(2)}$}
  %\nodetext{V32}(0.7,0){$w_2^{(2)}$}
  %
 %
 % EDGES
 \graphlinewidth{0.025}
% % First level
%  \edge{V0}{V11}
%  \edge{V0}{V12}

 % Second level
 \bow{V21}{V11}{0.09}
  \bow{V21}{V11}{-0.09}
    \edge{V22}{V11}
     
 \bow{V22}{V12}{0.09}
 \bow{V22}{V12}{-0.09}
 \edge{V22}{V12}
 
 \edge{V23}{V13}
 \bow{V23}{V13}{0.09}
 \bow{V23}{V13}{-0.09}
 \edge{V23}{V11}

 \edge{V24}{V14}
 \bow{V24}{V14}{0.09}
 \bow{V24}{V14}{-0.09}
 \edge{V24}{V11}
 

 %third level
 \bow{V31}{V21}{0.09}
 \bow{V31}{V21}{-0.09}
   \edge{V32}{V21}
%     \edge{V32}{V22}[\graphlinecolour(1,0,0)]
 \bow{V32}{V22}{0.09}
 \bow{V32}{V22}{-0.09}
 \edge{V32}{V22}
 
 

 \bow{V33}{V23}{0.09}
 \bow{V33}{V23}{-0.09}
 \edge{V33}{V21}
  \edge{V33}{V23}

 \bow{V34}{V24}{0.09}
 \bow{V34}{V24}{-0.09}
 \edge{V34}{V21}
  \edge{V34}{V24}  
 
    \freetext(8.9,5){$\ldots$}
  \freetext(8.9,3){$\ldots$}
    \freetext(8.9,1){$\ldots$}
    \freetext(2,0.1){$\vdots$}
  \freetext(4,0.1){$\vdots$}
    \freetext(6,0.1){$\vdots$}
    \freetext(8,0.1){$\vdots$}
%\freetext(3,0.1){.\,.\,.\,.\,.\,.\,.\,.\,.\,.\,.\,.\,.\,.\,.\,.\,.\,.\,.\,.\,..\,.\,.\,.\,.\,.\,.\,.}
\end{graph}
\caption{A stationary diagram with non-transitive tail-equivalence relation.}\label{Fig:nontransitnonloccomp}
\end{figure}
\end{example}



We generalize the Proposition~\ref{prop:toptransirred} to obtain the following theorem:

\begin{thm}\label{thm:classII}
    Every generalized Bratteli diagram is isomorphic to an irreducible generalized Bratteli diagram. 
\end{thm}


\begin{proof} 
    %We present the proof for generalized Bratteli diagrams. The same proof can be also used for standard Bratteli diagrams. 
    Let $B$ be a generalized Bratteli diagram. Then $X_B$ is a Polish space, hence, there exist a countable set of paths $\{x^{i}\}_{i = 0}^{\infty} \subset X_B$ which is dense in $X_B$. We show that there exists an enumeration of the vertices of $B$ such that for all $n$, the set $V_n$ is identified with $\N_0$, and all the paths $\{x^{i}\}_{i = 0}^{\infty}$ pass through the vertex number $j$ infinitely many times for every $j \in \N_0$.
%\tcr{add comment about isomorphism as enumeration of vertices.}
%Theorem \ref{thm:classIII} can be generalized to the case when
% $$
% X_B = \bigcup_{i = 0}^{\infty} \ov{\mathcal{O}(x^i)}.
% $$
%We show that we can enumerate the vertices of $V_n$ by numbers from $\N_0$ so that every path $x^i$ passes through every number $\N_0$ infinitely many times. 
We use the construction which is similar to the construction of Toeplitz sequences (see, e.g., \cite{Kurka2003} or \cite{Downarowicz2005}). Enumerate the vertices of $B$ so that the path $x^0$ passes through the vertex number $0$ on level $0$. Then we do not specify through which vertex  $x^0$ passes on level $1$. Then let $x^0$ pass  through vertex $0$ on level $2$ and through vertex $1$ on level $3$. We do not specify though which vertices $x^0$ passes on levels $4$ and $5$. Let $x^0$ pass through vertex $0$ on level $6$, through vertex $1$ on level $7$, and vertex $2$ on level $8$. 
%$0,1,2$ on levels $7,8,9$ 
We do not specify though which vertices $x^0$ passes on levels $9,10,11$, 
and so on. In other words, we consider sequence of vertices through which the path $x^0$ passes. Starting with the first vertex $s(x^0)$ and $n = 1$, we enumerate consecutive blocks of vertices of length $n$ by the numbers $0, \ldots, n-1$, and we leave the next block of vertices of length $n$ without forcing any enumeration. We show the procedure below. We denote by $*$ the vertices of $x^0$ without specified enumeration. The positions of the numbers in the sequence below correspond to the levels where we enumerate the vertices of $x^0$ by the given numbers:
\begin{equation}\label{x^0}
0*01**012***0123****\ldots
\end{equation}
Now we enumerate some of the vertices of the path $x^1$. We will enumerate only those vertices which belong to the levels, for which we didn't specify enumeration for $x^0$. In such a way, we enumerate the vertices of $x^0$ and $x^1$ on independent levels. Moreover, among the blocks of vertices, which are not enumerated for $x^0$, we keep every second block of vertices not enumerated for $x^1$.  
In other words, for $x^1$, we enumerate vertices only from the levels that correspond to the blocks of stars $*$ of odd lengths. In those blocks of odd lengths $2n-1$, we fill in only initial positions of lengths $n$ by the symbols $0, \ldots, n-1$. Hence, $x^1$ passes through vertex $0$ on level $1$, then through vertex $0$ on level $9$, vertex $1$ on level $10$, and so on. We show the enumerated vertices for $x^1$ on the corresponding levels with barred symbols:
$$
0\ov{0}01**012\ov{0}\ov{1}*0123****\ldots
$$
%In other words, we enumerate the vertices of $x^1$ which belong to every second interval with corresponding block $0,1,2,\ldots,n$, we leave more than enough space for the block to fit in). 
We don't use anymore in the construction the levels that corresponded to the blocks of stars of odd lengths in \eqref{x^0}, even if we didn't fully fill them in. For instance, we will never specify the number of vertex on level $11$ for any path $x^i$, where $i \in \N_0$. 
Now we enumerate in a similar way some of the vertices of the path $x^2$, which belong to the levels, where we haven't enumerated vertices of $x^0$ and $x^1$. Again, among the blocks of stars from \eqref{x^0} of even lengths, we leave every second block not enumerated. The numbers of vertices for $x^2$ on the corresponding levels are shown below by underlined symbols:
$$
0\ov{0}01\underline{0}*012\ov{0}\ov{1}*0123****\ldots
$$
%Note that in this construction, some positions will be never filled in, for instance, we don't specify the enumeration of the vertices on levels $5$ or $11$ for any path $x^i$, $i \in \N_0$. 
We repeat the procedure for each $x^i$, $i \in \N_0$. Thus, every path $x^i$ will pass through every vertex $j \in \N_0$ infinitely many times.
%\tcr{to be clarified, add an example similar to Example 3.10 but for countably many subdiagrams???}

Now we show that with the enumeration above, the diagram $B$ becomes irreducible.
Let $n \in \N_0$ and $i,j \in \N_0$. Consider any cylinder set $[y_0, \ldots, y_n]$ such that $r(y_n) = i$. 
%(in the case when $n = 0$, we just consider a cylinder set [i]).\tcr{introduce this notation earlier}
%Indeed, every cylinder set in $X_B$ contains at least one of the paths $x^i$.
Then there exists $k \in \N$ such that $x^k = (x^k_l)_{l=0}^{\infty} \in [y_0, \ldots, y_n]$. In other words, $x^k_l = y_l$ for $l = 0, \ldots, n$. 
Since $x^k$ passes through all vertices number $j$ infinitely many times for every $j \in \N_0$, there exists level $m > n+1$ such that $s(x^k_m) = j$. Thus, there is a finite path between a vertex $i$ on level $n$ and a vertex $j$ on some level $m > n$.
%\tcr{maybe separately n = 0}
\end{proof}

\begin{remark}
Note that Theorem \ref{thm:classII} also holds for standard Bratteli diagrams with the same number of vertices on each level, with a similar proof.
\end{remark}


%\tcr{compare to standard BD: what is possible}

Below we provide examples of completely and relatively irreducible generalized Bratteli diagrams, and give necessary conditions and sufficient conditions for a diagram to be completely or relatively irreducible.

\begin{example}[A completely irreducible generalized Bratteli diagram]\label{ex:GBDclass1} The Bratteli diagram $B_{RS}$ corresponding to the renewal shift (see Subsection \ref{subsec:BD_renewal}) is completely irreducible. Clearly, diagram $B_{RS}$ is irreducible: for every level $n$ and every vertex $i \in V_n$ we can reach vertex $0 \in V_{n+i}$, and from vertex $0  \in V_{n+i}$ we can reach any other vertex $j \in V_{n + i + 1}$. This proof does not depend on the enumeration of the vertices. In other words, let $B'=(V',E')$ be isomorphic to $B_{RS}$ and $g_n \colon V_n \rightarrow V_n'$, $n \geq 0$ be the corresponding maps from the definition of isomorphism. Then $B'$ has the unique vertex $g_n(0)$ on each level $n$ such that $E(g_n(0),v) \neq \emptyset$ for all $v \in V'_{n+1}$. In the Bratteli diagram $B_{RS}$, for all $n$, the vertex $0 \in V_{n+i}$ is joined by a path of length $i$ to the vertex $i \in V_n$. Hence, for all $n, i \in \N_0$ the vertex $g_{n+i}(0) \in V'_{n+i}$ is joined by a path of length $i$ to the vertex $g_n(i) \in V_n'$. 
Moreover, for every $n$ and every vertex  $i' \in V'_n$, there exists $i \in V_n$ such that $i' = g_n(i)$. Thus, for all $n$ and all $i' \in V'_n$, we have $E(i', g_{n+i}(0)) = E(g_n(i), g_{n+i}(0)) \neq \emptyset$. Therefore, for every $n$ and every vertex $i'\in V'_n$ we can reach the corresponding vertex $g_{n+i}(0)$, and from $g_{n+i}(0)$ we can reach any vertex $j$ on level $V'_{n+i+1}$. Hence, the diagram $B'$ is irreducible.
\end{example}

% \begin{prop}\label{Prop:renewal_irred_isom}
%     Let $B = (V,E)$ be the stationary generalized Bratteli diagram corresponding to the renewal shift. Then $B$ is irreducible \tcr{and connected as an undirected graph.} Moreover, for any $B' = (V',E')$ isomorphic to $B$, diagram $B'$ is also irreducible \tcr{and connected as an undirected graph.}
% \end{prop}

% \begin{proof}
        
%     \begin{figure}[hbt!]
% \unitlength=1cm
% \begin{graph}(9,6)
% % \graphnodesize{0.2}
% % \roundnode{V0}(3,6)
% %  %\nodetext{V0}(-1,0){$V_0$}
%  \roundnode{V11}(0,5)
%  %\nodetext{V11}(-0.7,0){$w_1^{(0)}$}
%   %\nodetext{V12}(0.7,0){$w_2^{(0)}$}
%  \roundnode{V12}(2,5)
%  \roundnode{V13}(4,5)
%  \roundnode{V14}(6,5)
%   \roundnode{V15}(8,5)


%   %
%     % The second level vertices
%    \roundnode{V21}(0,3)
%  %\nodetext{V11}(-0.7,0){$w_1^{(0)}$}
%   %\nodetext{V12}(0.7,0){$w_2^{(0)}$}
%  \roundnode{V22}(2,3)
%  \roundnode{V23}(4,3)
%  \roundnode{V24}(6,3)
%   \roundnode{V25}(8,3)
%   % The third level vertices
%  \roundnode{V31}(0,1)
%  \roundnode{V32}(2,1)
%  \roundnode{V33}(4,1)
%   \roundnode{V34}(6,1)
%     \roundnode{V35}(8,1)
  
%  %
%  % EDGES
%  \graphlinewidth{0.025}
% % % First level
% %  \edge{V0}{V11}
% %  \edge{V0}{V12}


%     \edge{V21}{V11}

%     \edge{V22}{V11}
%     \edge{V21}{V12}

    
     
%  \edge{V23}{V11}
%  \edge{V22}{V13}

    
%    \edge{V24}{V11}
%     \edge{V23}{V14}
    
    
%    \edge{V25}{V11}
%     \edge{V24}{V15}
    
%  % Second level

%     \edge{V31}{V21}

%     \edge{V32}{V21}
%     \edge{V31}{V22}

    
     
%  \edge{V33}{V21}
%  \edge{V32}{V23}

    
%    \edge{V34}{V21}
%     \edge{V33}{V24}

    
%    \edge{V35}{V21}
%     \edge{V34}{V25}
    
    
%   \freetext(8.9,5){$\ldots$}
%  \freetext(8.9,3){$\ldots$}
%   \freetext(8.9,1){$\ldots$}
  
%   \freetext(0,0.5){$\vdots$}
%   \freetext(2,0.5){$\vdots$}
%     \freetext(4,0.5){$\vdots$}
%       \freetext(6,0.5){$\vdots$}  
%         \freetext(8,0.5){$\vdots$} 



% \end{graph}
% \caption{}
% \end{figure}

% Notice that for every level $n$ there is a vertex $w_n \in V_n$ such that $w_n$ has outgoing edges to all vertices from $V_{n+1}$. This property is preserved under isomorphism of generalized Bratteli diagrams.
% Notice also that for every vertex $w$ on every level $n$ there is a finite path from from $w$ to $w_m$ for some $m > n$. This property is also preserved under isomorphism.
% Since from $w_m$ there is an edge to any vertex $v \in V_{m+1}$, we obtain that for any $n \in \mathbb{N}_0$, every $w' \in V'_{n}$ and every $v'$ there is an $m > n$ such that $E(w',v') \neq \emptyset$ for $v' \in V'_m$.
% \end{proof}



The following proposition is a straightforward generalization of Example~\ref{ex:GBDclass1}. 

\begin{prop}\label{Prop:suff_cond_irred_isom}
    Let $B$ be a generalized Bratteli diagram. Assume that
    
    (i) for every $n \geq 0$ there is a vertex $w_n \in V_n$ such that $E(w_n,v) \neq \emptyset$ for all $v \in V_{n+1}$,

    (ii) for every vertex $w$ on every level $n$ there is a finite path from from $w$ to $w_m$ for some $m > n$.

    Then $B$ is completely irreducible.
\end{prop}

Note that condition (i) in Proposition~\ref{Prop:suff_cond_irred_isom} is not sufficient for a diagram $B$ to be irreducible. Indeed in the diagram $B_{\infty}$ (see Subsection \ref{subsec:diagramB-infty}), the first vertex of each level has an edge to all vertices on a level below, but the diagram is reducible.



\begin{example}[Relatively irreducible Bratteli diagrams] Bratteli diagrams of bounded size (see Subsection \ref{subsec:bdd}) are relatively irreducible. 
Indeed, let $B=(V,E)$ be a generalized Bratteli diagram of bounded size  with parameters $(t_n, L_n)$. For every $w \in V_0$, let $Y_w^+$ denote the set of all infinite paths which start at $w$ and then pass through the rightmost possible vertex on each level, i.e. for every $m \in \mathbb{N}$, the paths from $Y_w^+$ go through the vertex $w + \sum_{i = 0}^{m-1} t_i$ on level $m$.

% We note that the isomorphism of generalized Bratteli diagrams \tcr{what about finite ones?} doesn't preserve irreducibility:
% \tcr{When is bdd size BD irreducible}
% \tcr{Formulate the result in terms of matrices}

\begin{prop}
    Let $B$ be a generalized Bratteli diagram of bounded size. Then $B$ is relatively irreducible. %isomorphic to a reducible generalized Bratteli diagram.
\end{prop}


\begin{proof}
    Let $(t_n,L_n)$ be the parameters for $B$ from the definition of a diagram of bounded size. For every $n \geq 1$ and $v \in V_n$, let $g_n(v) = v - \sum_{i = 0}^{n-1} t_{i}$ and $h_n$ be the corresponding bijection between the edges. Then in the new diagram $B'$, the images of the slanted paths from $Y_w^+ \subset X_B$ will be vertical paths, and any vertex $w \in V_0$ will be joined by finite paths only with vertices $v \leq w$ on lower levels.
\end{proof}

Clearly, all reducible generalized Bratteli diagrams such as diagrams of infinite odometers (see Subsection \ref{subsec:DIO}) or the diagram $B_{\infty}$ (see Subsection \ref{subsec:diagramB-infty}) are relatively irreducible.

\end{example}

%\begin{example}[Bratteli diagram of Class III] Bratteli diagram $B(2^n)$ belongs to Class III (or infinite Pascal from paper with Artem)?  \end{example} -- not true

%\subsection{Generalized Bratteli diagrams isomorphic only to irreducible Bratteli diagrams}



Throughout the paper, by a neighbourhood of a point $x$ we mean any set $U$ that contains an open set containing $x$.

\begin{thm}\label{Thm:isomredus}
    Let $B$ be a generalized Bratteli diagram. 
    If $B$ has a point with a compact neighbourhood then $B$ is relatively irreducible. 
    %If $B$ is completely irreducible then $X_B$ cannot have a point with a compact neighbourhood. 
    %The converse is not true: there exist reducible Bratteli diagrams such that every cylinder set is non-compact.
    
    %If $X_B$ has a point with a compact  neighbourhood then $B$ is isomorphic to a reducible generalized Bratteli diagram, hence $B$ cannot belong to Class I.
\end{thm}

\begin{proof} %Assume that $B$ has a point with a compact neighbourhood.
Enumerate the vertices $V_n$ of each level of $B$ by integers. Assume that $X_B$ possesses a point $x = (x_n)$ with a compact neighbourhood $U$. Since all cylinder sets form the base of topology and are closed, there exists a cylinder set $[(x_0, \ldots, x_n)]$ which contains $x$ and is compact. Hence, the set $[(x_0, \ldots, x_n)]$ can be represented as a standard Bratteli diagram, i.e. from the vertex $r(x_n) \in V_{n+1}$ one can reach only finitely many vertices %$W_{n+i} \subset V_{n+i}$
    from $V_{i}$
    for each $i > n + 1$.
    %and the same property holds for all the vertices one can reach from $r(x_n)$. 
    Thus, the paths that pass through $r(x_n)$ form a cone similarly to the cones for the diagrams of bounded size. Let $v_{n + 2} \in V_{n+2}$ be the leftmost vertex that can be reached from $r(x_n)$. Let $v_{n+3}$ be the leftmost vertex that can be reached from $v_{n+2}$ and so on. Hence, we obtain a sequence of vertices $\{v_{i}\}_{i = n+2}^{\infty}$, through which passes the leftmost path of the downwards directed cone starting at $r(x_n)$. Consider a shift $\tau_i \colon V_i \rightarrow V_i$ for each level $i > n + 1$ which shifts vertex $v_i$ so that its number becomes equal to the number of $r(x_n)$. Hence, the leftmost path becomes vertical and we cannot reach from $r(x_n) \in V_{n+1}$ any vertex with number less that $r(x_n)$ on any level below. Thus, diagram $B$ is isomorphic to a reducible diagram and cannot be completely irreducible.
\end{proof}

\begin{example}
    Every diagram of bounded size (see Subsection \ref{subsec:bdd}) has a locally compact path space. The diagram $B_{RS}$ which corresponds to the renewal shift (see Subsection \ref{subsec:BD_renewal}) is completely irreducible and does not have points with a compact neighbourhood.
    Indeed, every path $x \in X_{B_{RS}}$ passes through the first vertex infinitely many times. Since the first vertex of each level of $B_{RS}$ has infinitely many outgoing edges, every cylinder set which contains $x$ is not compact.
    
    %If for every level $n \in \N_0$ there is a vertex $v_n$ which reaches all other vertices, it doesn't mean that the diagram $B$ is irreducible.      If every vertex has infinitely many outgoing edges, i.e. there are no paths in $X_B$ with compact neighbourhoods, it doesn't mean that the diagram is irreducible.     Indeed, the diagram $B_{\infty}$ satisfies both of the abovementioned conditions but is reducible, the first vertex on each level reaches all vertices on the level below and every vertex has infinitely many outgoing edges. Thus, the condition of the Proposition \ref{prop:isomredus} is sufficient for the diagram to be isomorphic to a reducible Bratteli diagram, but it is not necessary.
    \end{example}

\begin{remark}
% In other words, if a diagram is isomorphic only to irreducible Bratteli diagrams, it cannot have paths with compact neighbourhouds, and we obtain a
Theorem~\ref{Thm:isomredus} provides a 
necessary condition for a Bratteli diagram to be completely irreducible: such diagram cannot have paths with compact neighbourhouds. However, this condition is not sufficient. For the diagram $B_{\infty}$ (see Subsection \ref{subsec:diagramB-infty}), every vertex has infinitely many outgoing edges, hence there are no points with compact neighbourhoods, but the diagram is reducible. 
\end{remark}


For standard Bratteli diagrams, there is a notion of a simple diagram: it's a diagram such that for every level $n$ there is a level $m > n$ such that $E(v,w) \neq \emptyset$ for all $v \in V_m$ and $w \in V_n$. %Clearly, the same definition cannot work for generalized Bratteli diagrams since every vertex can have only finitely many incoming edges.
It is well known that a standard Bratteli diagram has a minimal tail equivalence relation if and only if the diagram is simple 
(see, e.g., \cite[Theorem 2.5]{Putnam2010}). For generalized Bratteli diagrams, one cannot use straightforwardly the notion of minimality. Since every set $V_n$ is countably infinite, and every vertex of a generalized Bratteli diagram can have only finitely many incoming edges, it is not possible for a vertex to be connected to all vertices on some level above. The following theorem motivates viewing completely irreducible generalized Bratteli diagrams as analogues of simple standard Bratteli diagrams.

%The proof uses compactness and the result is not true for generalized Bratteli diagrams. \tcr{add def of simple and more to corresp subsection}



\begin{theorem}\label{thm:Class1_minimalR}
    Let $B$ be a generalized Bratteli diagram and $\mathcal{R}$ be the tail equivalence relation on $B$. If $B$ is completely irreducible then $\mathcal{R}$ is minimal.
\end{theorem}  

\begin{proof} Assume the contrary. Suppose $\mathcal{R}$ is not minimal. Then there exists a point $x = (x_n) \in X_B$ such that $X_B \setminus \ov{\mathcal{O}(x)} \neq \emptyset$. Hence there exists a cylinder set $[(e_0, \ldots, e_n)]$ which belongs to $X_B \setminus \ov{\mathcal{O}(x)}$. Since $[(e_0, \ldots, e_n)] \cap {\mathcal{O}(x)} = \emptyset$, we have $r(e_n) \neq r(x_n)$.
Consider a diagram $B' = (V', E')$ which is isomorphic to $B$ with the corresponding maps $g_m \colon V_m \rightarrow V_m'$ such that $g_n(r(e_n)) = 0 \in V'_{n+1}$ and $g_m(r(x_m)) = 1 \in V'_{m+1}$ for all $m \geq n$. 
Then there is no finite path between $0 \in V'_{n+1}$ and vertex $1 \in V'_{m+1}$ for all $m > n$. Indeed, if such a path existed, then there would exist a finite path between $r(e_n)$ and $r(x_m)$ in $B$, hence the orbit of $x$ would visit $[(e_0, \ldots, e_n)]$ and we would get a contradiction. Thus, $B'$ is reducible and $B$ cannot be completely irreducible.
%\tcr{SB: How do you know that such $(g_n)$ exists? Should you construct $B'$?}
\end{proof}

\textit{Open problem.} Assume that the tail equivalence relation $\mathcal{R}$ on a generalized diagram $B$ is minimal. Is it true that $B$ is completely irreducible?

\medskip
Proposition \ref{prop:minloccomp} might be a step towards solving the open problem above.

\begin{prop}\label{prop:minloccomp}
If tail equivalence relation is minimal then either $X_B$ is locally compact or there is no path in $X_B$ with a compact neighbourhood.    
\end{prop}

\begin{proof}
    Assume that there is a path $x = (x_n) \in X_B$ without a compact neighbourhood. Then for all $n$, the vertex $s(x_{n})$ is joined by a path to a vertex on some level below which has infinitely many outgoing edges. 
    Since $\mathcal{R}$ is minimal, the orbit $\mathcal{O}(x)$ is dense in $X_B$. Assume that there is a point $y = (y_n) \in X_B$ with a compact neighbourhood. Then there exists $N \in \mathbb{N}$ such that the cylinder set $[(y_0, \ldots, y_N)]$ is compact. Since $\ov{\mathcal{O}(x)} = X_B$, there is a finite path between $r(y_N)$ and $s(x_m)$ for some $m \geq N$. Hence, there is a path between $r(y_N)$ and a vertex which has infinitely many outgoing edges. Hence, $[(y_0, \ldots, y_N)]$ cannot be compact. 
    %$s(x_k)$ for all $k \geq m$ and thus there is a path between $r(y_N)$ and $s(x_{n_i})$ for $n_i \geq m$. It follows that $[(y_0, \ldots, y_N)]$ is not compact and we get a contradiction.
\end{proof}




%\tcr{double check the proof and what we use}
%\tcr{is the notion of minimality Ok in Polish space}
%\tcr{other propeties like mixing?}






%\tcr{result below: what's known for standard bd. Ask Artem? For fin rank bd?}

%\tcr{Hypo: if there are uncountably many disjoint orbit closures, i.e. when one can embed a dyadic-odometer binary tree into GBD as a cylinder set, then GBD belongs to Class III.}

% \tcr{top transitive point into top transitive?}
% \tcr{prove everything and make it one theorem}
% \tcr{not top transitive <=> splits into disjoint subdiagrams after some level?}



% \begin{theorem}
%     Minimality and topological transitivity are preserved under isomorphism.
% \end{theorem}

% \begin{proof}
%     Assume that $B$ is a generalized Bratteli diagram such that the tail equivalence relation $\mathcal{R}$ is topologically transitive. Hence there is a path $x = (x_n) \in X_B$ such that the orbit of $x$ under the tail equivalence relation visits every cylinder set. Thus, for every $n \geq 0$ and for every cylinder set $[(y_0, \ldots, y_n)]$, there is a finite path between $r(y_n)$ and $s(x_m)$ for some $m > n$. Clearly, there exists a path between $g(r(y_n))$ and $g(s(x_m))$ which is an image of the path between $r(y_n)$ and $s(x_m)$ under the isomorphism (the renumbering of vertices does not change the fact that there is a path between them). So the image of $x$ under the isomorphism will have a dense orbit in the path space of the image of $B$ under the isomorphism. Similarly, minimality of the tail equivalence relation is preserved under isomorphism. \tcr{Are all ``top'' properties preserved?}
% \end{proof}


% \subsection{Two-sided diagram of infinite odometers}

% Unlike the case of $B_{IO}$ considered in \cite{BezuglyiKarpelKwiatkowski2024}, the stationary two-sided $\tl B_{IO}$ can have infinitely countably many probability ergodic invariant measures. 

% Horizontally stationary and vertically stationary 2-sided DIO has no probability ergodic invariant measure.

% All bounded size Bratteli diagrams (vertices are enumerated by $\mathbb{Z}$) can be made reducible after the change of enumeration of the vertices (graph isomorphism).

% \subsection{Tail equivalence relation}

% Notice that even in the case of a stationary generalized Bratteli diagram the existence of a vertex that has an outgoing edge to every vertex on the level below doesn't guarantee that the tail equivalence relation is minimal (see Example~\ref{Ex:ReducibleGBD1VertexAccessToAll}).

% \begin{example}\label{Ex:ReducibleGBD1VertexAccessToAll}
% Let $B$ be a stationary generalized Bratteli diagram with the set of vertices on each level identified with $\mathbb{N}$ and the entries of the incidence matrix given by the formula 
% $$
% f_{ij} = 
% \left\{
% \begin{aligned}
% & 1, \mbox{ for } i \geq j,\\
% & 0, \mbox{ otherwise. }
% \end{aligned}
% \right.
% $$
% Then we have
%     $$F = 
%     \begin{pmatrix}
%      1 & 0 & 0 & 0 &\ldots\\
%      1 & 1 & 0 & 0 &\ldots\\
%      1 & 1 & 1  & 0 &\ldots\\
%      1 & 1 & 1  & 1 & \ldots\\
%      \vdots & \vdots & \vdots & \vdots & \ddots
%     \end{pmatrix},
%     $$
% matrix $F$ has an infinite sum of elements in every column and the finite sum of elements in every row. For every $n \in \mathbb{N}_0$ and every $i \in V_n$, there is an edge from $i$ to the vertices $i, i + 1, i +2, \ldots \in V_{n+1}$. In particular, the first vertex of each level has outgoing edges to every vertex on the next level. It is easy to see that no vertical path is topologically transitive. Moreover, the path that passes through the first vertex of each level is a fixed point for the tail equivalence relation. A slanted path $y = (y_n)$ which passes through the vertices $s(y_n) = n \in V_n$ is a topologically transitive point for $\mathcal{R}$. It follows from the fact that for every $n$, from the vertex $n \in V_n$ we can reach by a finite path any vertex $i \in V_m$ for $i \leq n$ and $m < n$. Thus, with growing $n$, we will be able eventually to reach from the vertices of $s(y_n)$ any vertex in $B$.

% \tcr{I suppose that the matrix is transient and there exists
% a finite invariant measure.}

% \begin{figure}[hbt!]
% \unitlength=1cm
% \begin{graph}(9,4)
%   %
%     % The second level vertices
%    \roundnode{V21}(0,3)
%  %\nodetext{V11}(-0.7,0){$w_1^{(0)}$}
%   %\nodetext{V12}(0.7,0){$w_2^{(0)}$}
%  \roundnode{V22}(2,3)
%  \roundnode{V23}(4,3)
%  \roundnode{V24}(6,3)
%   \roundnode{V25}(8,3)
%   % The third level vertices
%  \roundnode{V31}(0,1)
%  \roundnode{V32}(2,1)
%  \roundnode{V33}(4,1)
%   \roundnode{V34}(6,1)
%     \roundnode{V35}(8,1)
  
%  %
%  % EDGES
%  \graphlinewidth{0.025}
% % % First level
% %  \edge{V0}{V11}
% %  \edge{V0}{V12}



    
%  % Second level

%     \edge{V31}{V21}
    
%     \edge{V32}{V21}
%         \edge{V32}{V22}
     
%  \edge{V33}{V21}
%   \edge{V33}{V22}
%  \edge{V33}{V23}
    
%    \edge{V34}{V21} 
%     \edge{V34}{V22} 
%        \edge{V34}{V23} 
%    \edge{V34}{V24} 
   
%    \edge{V35}{V21}
%       \edge{V35}{V22}
%          \edge{V35}{V23}
%             \edge{V35}{V24}
%    \edge{V35}{V25}
    
    
%  \freetext(8.9,3){$\ldots$}
%   \freetext(8.9,1){$\ldots$}
  
%   \freetext(0,0.5){$\vdots$}
%   \freetext(2,0.5){$\vdots$}
%     \freetext(4,0.5){$\vdots$}
%       \freetext(6,0.5){$\vdots$}  
%         \freetext(8,0.5){$\vdots$} 


% \end{graph}
% \caption{}\label{Fig:ReducibleGBD}
% \end{figure}

% \end{example}

\subsection{Connectedness of Bratteli diagrams}
In this subsection, we discuss connectedness of Bratteli diagrams as undirected graphs. Recall that a graph is called \textit{connected} if there exists a path between every pair of vertices. We discuss the relation between connectedness of a generalized Bratteli diagram, ireducibility, and topological properties of the tail equivalence relation.

%\tcr{no sense for directed BD} or weakly connected directed graph

%In this section, we study the notion of isomorphism for generalized Bratteli diagrams, and find out which properties of Bratteli diagrams are preserved under isomorphism. For instance, in general, irreducibility is not preserved under isomorphism (see \cite[Remark 2.10]{BezuglyiJorgensenKarpelSanadhya2025}). We investigate for which classes of (generalized) Bratteli diagrams irreducibility is preserved under isomorphism. \tcr{Also consider for standard BD?} Notice that in the case of standard Bratteli diagrams, the notion of simplicity is usually used instead of irreducibility, and simplicity is preserved under isomorphism. In the case of generalized Bratteli diagrams, we cannot consider the notion analogous to that of a simple Bratteli diagram since every vertex of a generalized Bratteli diagram has only finitely many incoming edges. Thus, by definition, a generalized Bratteli diagram cannot be simple. 

\begin{remark}
    Note that for standard Bratteli diagrams, it is usually considered that the level $V_0$ consists of a single vertex $v_0$ which is connected to all vertices on level $V_1$. Thus, every vertex $v$ on any level $V_n$, for $n > 0$, is connected by a finite path to $v_0$. With such definition, standard Bratteli diagrams
    are always connected.
\end{remark}

The example below shows that, in general, the notions of irreducibility and connectedness for a generalized Bratteli diagram are independent. 

\begin{example} Diagram $B_{RS}$ (see Subsection \ref{subsec:BD_renewal}) is irreducible and connected. Diagram from Example \ref{ex:bdd_irred_nontrans} is irreducible but not connected. Diagram $B_{\infty}$ (see Subsection \ref{subsec:diagramB-infty}) is connected but not irreducible. Diagram from Example \ref{ex:bdd_red_nontrans} is reducible and not connected.
\end{example}

% \begin{example}
%     Diagram $B(2^n)$ is not connected while diagram for renewal shift is connected.
%     \tcr{generalize}
% \end{example}

Proposition \ref{prop:Rconnect} demostrates the relation between connectedness of a generalized Bratteli diagram and topological transitivity of the corresponding tail equivalence relation.

\begin{prop}\label{prop:Rconnect}
    Let $B$ be a generalized Bratteli diagram and $\mathcal{R}$ be the tail equivalence relation on $B$. If $\mathcal{R}$ is topologically transitive then $B$ is connected.
\end{prop}

\begin{proof}
    Let $m,n \in \mathbb{N}_0$, $i \in V_m$ and $j \in V_n$. We show that there exists a finite path in the unordered graph $B$ which joins $i$ and $j$. Let $x = (x_k) \in X_B$ be a path such that $\ov{\mathcal{R}(x)} = X_B$. Then there exist paths $y  = (y_k), z  = (z_k) \in X_B$ such that $y, z \in \mathcal{R}(x)$ and path $y$ visits a cylinder that ends in vertex $i$, path $z$ visits a cylinder that ends in vertex $j$. In other words, $s(y_m) = i$ and $s(z_n) = j$. Since $y, z \in \mathcal{R}(x)$, there exist levels $M, N \in \mathbb{N}$ such that $y_k = x_k$ for all $k \geq M$ and $z_k = x_k$ for all $k \geq N$. Then a portion of path $y$ provides a finite path between vertex $i = s(y_m)$ and vertex $s(y_M) = s(x_M)$, a portion of path $x$ connects $s(x_M)$ and $s(x_N)$ and a portion of path $z$ connects $s(x_N) = s(z_N)$ and $s(z_n) = j$.
\end{proof}



\begin{corol}
    Let $B$ be an irreducible stationary generalized Bratteli diagram. Then $B$ is connected.
\end{corol}

\begin{proof}
    The proof follows directly from Proposition \ref{prop:Rconnect} and Theorem \ref{ThmBJKStoptrans}.%\cite[Theorem 5.1]{BezuglyiJorgensenKarpelSanadhya2025}.
\end{proof}

\begin{remark}
    Note that the converse to Proposition \ref{prop:Rconnect} is not true. Indeed, diagram form Example \ref{ex:connected_non-transR} is connected, but its tail equivalence relation is not topologically transitive.
\end{remark}

%Note that there are reducible diagrams of bounded size (see \cite{})





\textbf{Acknowledgements.} The author is very grateful to the colleagues and collaborators, 
especially, S. Bezuglyi, H. Bruin, P. Jorgensen, J. Kwiatkowski, P. Oprocha, S. Radinger, S.~Sanadhya, T. Raszeja  for valuable discussions.
This research was partially supported by a subsidy from the Polish Ministry of Science and Higher Education for the AGH University of Krakow.

\bibliographystyle{alpha}
\bibliography{ReferencesGBD2}

\end{document}
